% Modernised reading — fonds Alexandre Grothendieck, shelfmark 142, pages 1-121.
%
% This is an INTERPRETATION, not a transcription. It re-reads the pages in
% today's notation and vocabulary, states the hypotheses the manuscript leaves
% implicit, and drops the critical apparatus so that the mathematics reads
% straight through. Everything the brackets and underlines carried moves into a
% footnote. It is held to being mathematically correct as it stands; where the
% manuscript is loose or mistaken the true statement appears and the footnote
% says what the page has. In any dispute of substance the transcription
% governs: whoever cites Grothendieck cites batch-01.fr.tex ... batch-07.fr.tex.
%
% Scope: the folder is 121 pages in seven batches (1-20, 21-40, 41-60, 61-80,
% 81-100, 101-120, 121), all transcribed, read whole before any of this was
% written, and written as ONE file for the whole shelfmark. Blank or
% count-only sheets: 22, 42, 64; covers in his hand: 1, 21, 65, 104; pages 106
% and 108 have no transcription (versos without text, by the page numbering).
%
% The spine. Five runs, each with its own title or start, kept distinct:
%   I.   pp. 1-20  « Catégories de Galois-Teichmüller d'un corps alg. clos »
%        (his pp. 1-10): isotopy category of surfaces with finite étale maps,
%        its profinite completion, the Teichmüller-Galois category, the
%        hole-filling problem.
%   II.  pp. 21-41 « Opérations de Gamma_{Qbar/Q} sur pi_1 M_{0,3} » (his
%        pp. 5-23; his pp. 1-4 are not in the folder): seven groupoids I-VII on
%        the thrice-punctured sphere and the automorphisms commuting with the
%        symmetries; the « relation des lacets ».
%   III. pp. 43-63 « La sphère à trois points marqués » (his pp. 1-12): an
%        invariant set-up of the same sphere, real blow-up, base points and
%        base regions, « épinglages ».
%   IV.  pp. 65-103 « Généralités sur les actions de groupes de Galois sur les
%        compactifiés profinis de groupoïdes fondamentaux » (his pp. 1-18 and
%        a tail): cocycles, base regions T, Mordell-Weil, reduction to U_{0,3}
%        via Belyi, a lifting lemma for groupoid coverings, necessary
%        conditions for an automorphism to come from Galois.
%   V.   pp. 104-121 « Groupoïde fondamental d'une carte trouée (préliminaires
%        à l'action de Galois) »: sketches (105-112), then a run (his pp. 1-7,
%        113-121): presentations of the fundamental groupoid of a punctured
%        map, and hole-filling made trivial.
% II and III use the same letters (R, Q, P, a, b, c, d, phi, psi) for
% different arcs; V reuses tau_i, rho_i for the cartographic group. Each is
% read in its own notation and the spine section says so.
%
% Notation fixed for the whole folder. Paths compose as functions (right to
% left); for a path c : x -> y, c(h) = c h c^-1 transports pi(x) to pi(y).
% Pi_{0,3} = pi_1(P^1(C) - {0,1,inf}) = < l_0, l_1, l_inf | l_inf l_1 l_0 = 1 >,
% the order of the relation being the one of II and III (and of 145's second
% draft); Pi^_{0,3} its profinite completion. rho, sigma_inf, sigma~_inf,
% tau_inf are the automorphisms of Pi_{0,3} of pages 23 and 53-54. T_X =
% A_X / A°_X is « le groupe de Teichmüller » of the pages, today the extended
% mapping class group; T_X^+ the oriented one. The Galois group of IV is
% written Gamma_K where the page writes Gamma.
%
% Substantive departures from the pages, each footnoted where it occurs:
%   - p. 6: the proof of Prop. 2 is completed by a Lefschetz argument
%     (ours); the page only sketches it and calls the end « bien connu ».
%   - p. 7, Cor. 2: needs the contractibility of A°_X (Hamstrom), not stated.
%   - p. 11: the identification (17) of the two contracted products is
%     asserted, not proved; it needs a compatibility of profinite topologies
%     that the page does not discuss.
%   - pp. 11-14: reading order 11 -> 13 and 12 -> 14 (the text says so; the
%     transcription records the leaves in scan order).
%   - p. 14: « pas besoin qu'il soit de car. 0 » is restricted: in
%     characteristic p the fundamental group of an affine curve is not
%     invariant under extension of algebraically closed fields.
%   - p. 46: the real blow-up of the sphere at three points has 3
%     cross-caps; « genre 2 » is its first rational Betti number.
%   - p. 51: the incisions of Sigma~_eps are the d^{eps'} (as drawn), and the
%     arcs b_j are equatorial (the page writes supp c_j).
%   - p. 54: tau_inf(l_inf) = l_0 l_1 (page 23), not l_0^-1 l_inf^-1 l_0;
%     Pi~_{0,3} is PGL_2(Z) (as page 53 writes), not SL_2(Z).
%   - p. 55: with Sigma oriented, 3 resp. 2 épinglages, not « trois resp. ».
%   - pp. 35-36, (60)/(63): alpha' = zeta' rho(zeta')^-1, as the derivation of
%     p. 35 and (78) of p. 40 give, not rho^-1.
%   - p. 74, (26): the transition iso goes pi_alpha -> pi_beta.
%   - pp. 75-76: the subdivision argument needs K-points everywhere; the
%     correct statement is Brown's groupoid van Kampen theorem, which is the
%     page's own « résultat un peu plus général » with I_alpha = I cap T_alpha.
%   - p. 86: Belyi needs the curve defined over a number field.
%   - p. 114, (9): a^i_r : R^{i+1,i+2}_r -> R^{i+2,i+1}_r (the arc along the
%     side opposite i, as (10) requires); then (12) is the boundary of the
%     truncated triangle, a loop at R^{01}_r.
%   - p. 120, (21): A^{01}_r = (a^inf_{tau_0 r})^-1 Lambda'^1_r a^inf_r (the
%     page's inverse on Lambda'^1_r does not compose).
% The computations of II (pp. 23-41) depend on conventions fixed in his
% missing pp. 1-4; they are reported as the pages give them. The implications
% between equations that do not depend on those conventions — (23) => (12a),
% (37) => (26), (30), (32), (61), (63) <=> (70'), the derivation of (60) —
% were checked by hand.
%
% Announced and never established in the folder: the identification (17);
% an algebraic determination of the Teichmüller-Galois category (« programme
% à l'horizon », p. 16); a category containing both étale maps and
% hole-filling (pp. 17-20); the « formules explicites » for f_l(u) (p. 81);
% the conjecture that the conditions a), b) / 1)-4) are sufficient (p. 95);
% a conjecture « en termes de M_{0,4}, M_{0,5} » (p. 99, mostly illegible);
% the condition 5° (p. 101).
%
% Pass: Opus 5.5 (claude-opus-5-5), 2026-10-10 — first pass, unchecked against the pages by a human.
\documentclass[11pt,a4paper]{article}
\input{../preamble/grothendieck}

\folder{142}
\pages{1}{121}
\dating{[à partir de 1978]}
\watermark{Édition de démonstration\\interprétation personnelle de l'œuvre}
\foldertitle{Action de Galois sur Teichmüller : notes manuscrites (s.d.) — lecture modernisée du dossier entier}

\begin{document}

\section*{Résumé}

\begin{resume}
Ces cent vingt pages cherchent à voir le groupe de Galois des nombres
algébriques agir sur des objets qu'on peut dessiner : des surfaces, des
lacets tracés dessus, des cartes.

Le point de départ est un fait que tout le dossier exploite. Une surface
compacte dont on a percé quelques trous peut être regardée de deux façons :
comme un objet de topologie, qu'on peut déformer à volonté, ou comme une
courbe algébrique, définie par des équations. La première façon ne connaît que
la forme ; la seconde connaît aussi les nombres qui entrent dans les
équations. Or certains objets topologiques — les revêtements finis de la
surface, les lacets « vus de loin », c'est-à-dire à travers tous les
revêtements finis à la fois — se laissent décrire de la seconde façon. Dès
lors, toute symétrie des nombres, tout élément du groupe de Galois, les
transforme. On ne voit pas ces transformations : elles ne viennent d'aucune
déformation de la surface. Mais on peut chercher à les attraper par les
traces qu'elles laissent.

Les pages 1 à 20 posent le décor le plus large. On assemble toutes les
surfaces trouées, de tous les genres, en une seule catégorie, dont les flèches
sont les revêtements finis pris « à déformation près » ; on la complète, en y
ajoutant des limites, comme on passe des entiers aux entiers $p$-adiques ; et
l'on obtient un objet sur lequel le groupe de Galois agit, la \emph{catégorie
de Teichmüller–Galois}. Déterminer cet objet avec son action, par des moyens
purement combinatoires, c'est le « programme à l'horizon ». Un obstacle reste
en suspens : l'opération qui bouche un trou ne trouve pas sa place dans ce
cadre.

Le reste du dossier redescend au cas le plus petit : la sphère privée de trois
points. Pour y décrire l'action de Galois sur les lacets, il faut des points
de départ ; et c'est là que se joue une difficulté de méthode que Grothendieck
affronte page après page. Un point de départ unique brise toutes les
symétries de la situation : on fait les comptes depuis un carrefour choisi
arbitrairement, et chaque formule porte la trace de ce choix. Il remplace
donc le point par un \emph{système} de points de départ placés
symétriquement — tout près de chaque trou, aux pôles, aux points opposés —,
puis par des \emph{régions} de départ, des morceaux de sphère d'un seul
tenant, qu'on peut utiliser comme un point puisqu'on peut s'y promener sans
ambiguïté. Comme une ville où l'on se donne rendez-vous « sur la place » plutôt
qu'à une adresse : la place est assez grande pour que tout le monde s'y
retrouve, assez simple pour qu'il n'y ait qu'une façon d'aller d'un bout à
l'autre. Les symétries de la sphère déplacent ces lieux de rendez-vous les uns
sur les autres, et les calculs deviennent invariants. Il le dit lui-même :
« un aspect essentiel du travail à faire est de garder cette opération
constamment présente à l'esprit, et d'introduire des notations et des
constructions qui soient invariantes ». Et plus loin, se reprenant : « suivant
sans doute d'anciennes habitudes, j'ai privilégié » tels points de départ,
alors que ce sont d'autres qui expriment directement les phénomènes en jeu.

Ce que donnent ces calculs, ce sont des équations : une symétrie cachée des
lacets est codée par quelques éléments du groupe des lacets complété et par un
« multiplicateur », et ces éléments doivent satisfaire une « relation des
lacets ». Dix ans plus tard, ces équations deviendront, sous la plume de
Drinfeld et d'Ihara, les relations qui définissent le \emph{groupe de
Grothendieck–Teichmüller}. Les pages 65 à 103 en donnent le cadre général :
comment l'action d'un élément de Galois se lit par des « cocycles » attachés
aux chemins ; pourquoi, sur les courbes d'un certain type, un point rationnel
est entièrement déterminé par la façon dont Galois agit sur les chemins qui y
mènent ; pourquoi, grâce au théorème de Belyi, tout se ramène à la sphère à
trois trous ; et quelles conditions une symétrie des lacets doit remplir pour
venir de Galois — la première, et « de très loin » la plus importante selon
lui, étant précisément la relation des lacets. Il conjecture que ces
conditions nécessaires sont aussi suffisantes ; la question est aujourd'hui
encore ouverte.

Les dernières pages reviennent au concret : une carte dessinée sur une
surface, avec ses sommets, ses arêtes et ses faces percés, et les façons de
présenter ses lacets par un nombre fini de chemins et de relations. Là, la
difficulté laissée en suspens à la page 20 se dissout : boucher un trou, c'est
seulement écrire que le petit lacet qui en fait le tour devient trivial.

On lit dans ces notes une manière de travailler : une phrase qui s'arrête
parce que la notation ne convient pas, une notation reprise jusqu'à ce que les
symétries y soient visibles, un « travail fascinant en perspective, que je
n'ai pratiquement pas encore commencé ! ». La mention du théorème de Belyi
(1979) place au moins la quatrième suite après cette date ; le dossier est
voisin, dans l'inventaire, de \emph{La Longue Marche à travers la théorie de
Galois} (1981), dont il semble préparer plusieurs chapitres.

\keywords{mapping class group, Teichmüller groupoid, profinite completion of a category, Galois–Teichmüller theory, Grothendieck–Teichmüller group, real oriented blow-up, tangential base point, étale fundamental groupoid, non-abelian Galois cohomology, Kummer map, anabelian geometry, Belyi's theorem, van Kampen theorem for groupoids, covering of groupoids, dessins d'enfants}
\end{resume}

\pagerange{1}{121}
\section*{Le fil du dossier, et les conventions}

\subsection*{Les stations}

Le dossier réunit cinq suites, chacune avec son titre de couverture ou son
départ propre ; on les garde distinctes.

\begin{itemize}
\item \textbf{I. Catégories de Galois–Teichmüller (pages 1 à 20, ses pages 1
à 10).} La catégorie isotopique des surfaces trouées et des applications
finies étales ; sa complétion profinie, d'abord pour les surfaces, puis pour
une catégorie quelconque à groupes d'automorphismes profinis ; la catégorie de
Teichmüller–Galois d'un corps algébriquement clos et le « programme » ; le
problème du bouchage des trous, laissé ouvert. Les feuillets sont dans le
désordre : le texte court de la page 3 à la page 5, puis à la page 4, puis à
la page 6 ; et, plus loin, de la page 11 à la page 13, et de la page 12 à la
page 14.
\item \textbf{II. Opérations de $\Gamma_{\overline{\mathbb{Q}}/\mathbb{Q}}$
sur $\pi_1 M_{0,3}$ (pages 21 à 41, ses pages 5 à 23).} Ses pages 1 à 4
manquent au dossier, et avec elles les conventions de départ. Sept
groupoïdes $\mathrm{I}$ à $\mathrm{VII}$, chacun le groupoïde fondamental de
la sphère privée de trois points sur un système symétrique de points base ;
pour chacun, la description des automorphismes (du complété profini) qui
fixent les sommets et commutent aux symétries ; la « relation des lacets ».
\item \textbf{III. La sphère à trois points marqués (pages 43 à 63, ses pages
1 à 12).} La même sphère, reprise de façon invariante : géométrie euclidienne,
éclatement réel, points base sur les bords, régions simplement connexes,
isomorphismes d'« épinglage ».
\item \textbf{IV. Généralités sur les actions de groupes de Galois (pages 65 à
103, ses pages 1 à 18 et une suite).} Le formalisme des cocycles $f_l$,
$g_l$ ; les régions de base $T$ ; Mordell–Weil ; la réduction à $U_{0,3}$
par le théorème de Belyi ; un lemme de relèvement pour les revêtements de
groupoïdes ; les conditions nécessaires pour qu'un automorphisme provienne de
Galois.
\item \textbf{V. Groupoïde fondamental d'une carte trouée (pages 104 à 121).}
Des pages d'esquisses (105 à 112), puis une suite paginée de sa main de 1 à 7
(pages 113 à 121) : présentations du groupoïde fondamental d'une carte dont
sommets, milieux d'arêtes et centres de faces sont percés, et le bouchage des
trous.
\end{itemize}

L'ordre de rédaction de II et III n'est pas sûr. III installe des notations
qui semblent préparer II (« je vais réparer cette négligence » en privilégiant
les points $R_i^{\omega}$ voisins des trous, que II utilise d'emblée), mais
II a ses propres pages 1 à 4, perdues, et les lettres ne s'y accordent pas
entièrement. \textbf{Même nom, objets différents} : en II, $c_i$ va d'un
point $Q_i$ de l'équateur à un point voisin d'un trou et $d_i$ d'un pôle à
$Q_i$ ; en III, $c_i$ va d'un pôle à $Q_i$ et $d_i$ d'un pôle à un point
$S_i$ voisin du trou. Les $\varphi$, $\psi$ changent de rôle de même. En V,
$\tau_0, \tau_1, \tau_\infty$ et $\rho_0, \rho_1, \rho_\infty$ sont les
générateurs du groupe cartographique agissant sur les drapeaux d'une carte,
sans rapport avec les automorphismes $\tau_\infty$, $\rho$ de
$\Pi_{0,3}$ de II et III. Et $\varepsilon$ désigne en III à la fois un
hémisphère ($\varepsilon \in \boldsymbol{\varepsilon}$) et, avec un indice,
un élément de $\widetilde\Pi_{0,3}$ ($\varepsilon_0, \varepsilon_1,
\varepsilon_\infty$), comme sur les pages.

\subsection*{Les conventions}

\emph{Chemins.} Les chemins se composent comme des fonctions, de droite à
gauche : pour $l : x \to y$ et $l' : y \to z$, le composé est $l' l : x \to z$.
Pour un chemin $c : x \to y$, on note $c(h) = c\,h\,c^{-1}$ le transport
$\pi(x) \to \pi(y)$ qu'il induit sur les groupes fondamentaux. C'est la
convention des pages ; elle impose l'ordre de la relation ci-dessous.

\emph{Le groupe de la sphère à trois trous.} $U_{0,3} =
\mathbb{P}^1(\mathbb{C}) \setminus \{0, 1, \infty\}$,
\[
\Pi_{0,3} = \pi_1(U_{0,3}) = \langle\, l_0, l_1, l_\infty \mid l_\infty l_1 l_0 = 1 \,\rangle ,
\]
libre sur $l_0, l_1$ ; $\widehat\Pi_{0,3}$ est son complété profini. C'est
l'ordre de la seconde rédaction du dossier 145 (où ce groupe est noté $\pi$,
$\widehat\pi$)\footnote{Le dossier 145 souligne : « et non
$l_0 l_1 l_\infty = 1$ ! » ; sa première rédaction utilisait l'ordre
opposé.}. Les pages 23 et 53–54 fixent un formulaire d'automorphismes de
$\Pi_{0,3}$ :
\[
\rho : l_0 \mapsto l_1 \mapsto l_\infty \mapsto l_0, \qquad
\sigma_\infty : l_0 \leftrightarrow l_1,\ l_\infty \mapsto l_1^{-1} l_0^{-1} = l_0\, l_\infty\, l_0^{-1},
\]
\[
\tilde\sigma_\infty : l_0 \mapsto l_1^{-1},\ l_1 \mapsto l_0^{-1},\ l_\infty \mapsto l_\infty^{-1}, \qquad
\tau_\infty = \sigma_\infty\tilde\sigma_\infty = \tilde\sigma_\infty\sigma_\infty : l_0 \mapsto l_0^{-1},\ l_1 \mapsto l_1^{-1},\ l_\infty \mapsto l_0\, l_1 .
\]
On a vérifié que chacune de ces règles définit bien un automorphisme (la
relation $l_\infty l_1 l_0 = 1$ est respectée) et que $\sigma_\infty$ et
$\tilde\sigma_\infty$ commutent. $\rho$ est d'ordre $3$, $\sigma_\infty$,
$\tilde\sigma_\infty$, $\tau_\infty$ d'ordre $2$. Géométriquement, $\rho$
est la rotation d'un tiers de tour qui permute $0 \to 1 \to \infty$,
$\sigma_\infty$ le demi-tour qui échange $0$ et $1$, $\tau_\infty$ la
conjugaison complexe et $\tilde\sigma_\infty$ la réflexion qui échange $0$ et
$1$, chacune accompagnée d'un chemin qui ramène le point base à lui-même.

\emph{Teichmüller.} Pour une surface $X$, $A_X$ est le groupe de ses
homéomorphismes, $A^{\circ}_X$ le sous-groupe de ceux qui sont isotopes à
l'identité, et $T_X = A_X/A^{\circ}_X$ ce que les pages appellent le
« groupe de Teichmüller (non orienté) » — aujourd'hui le groupe modulaire
étendu, ou \emph{mapping class group} étendu ; $T^{+}_X$ est le sous-groupe
d'indice $2$ des classes qui respectent l'orientation\footnote{L'usage de
« groupe de Teichmüller » pour le groupe des classes d'isotopie est celui de
Grothendieck, ici comme dans l'\emph{Esquisse d'un programme} (1984) ;
l'usage courant le réserve au groupe modulaire agissant sur l'espace de
Teichmüller, ce qui revient au même pour les surfaces de type fini
anabéliennes.}. Une surface est \emph{anabélienne} si chacune de ses
composantes, de genre $g$ et à $\nu$ trous, vérifie $2g - 2 + \nu > 0$,
c'est-à-dire $\chi < 0$ — aujourd'hui : \emph{hyperbolique}.

\emph{Galois.} En IV, le groupe que la page note $\Gamma$ est noté
$\Gamma_K = \mathrm{Aut}(\mathbb{C}/K)$, pour le distinguer du groupe des
symétries de la sphère, noté $\Gamma_{\Sigma,J}$ en III.

\pagerange{1}{20}
\section*{I. Catégories de Galois–Teichmüller d'un corps algébriquement clos (pages 1 à 20)}

La page 1 est une couverture de sa main, qui donne ce titre, avec au crayon
« 19 pages ». La suite, paginée de 1 à 10, est un seul essai.

\pagerange{2}{8}
\subsection*{La catégorie isotopique (pages 2 à 8)}

Soient $X$, $Y$ deux surfaces topologiques orientées de type fini — des
surfaces compactes privées d'un nombre fini de points —, connexes pour
simplifier. Une \emph{application finie étale} $X \to Y$ est un revêtement
fini. On pose
\[
(1)\qquad \mathrm{Hom}_{\mathrm{isot}}(X, Y) = \pi_0\bigl(\text{espace des applications finies étales } X \to Y\bigr),
\]
l'ensemble des composantes connexes par arcs ; deux applications dans la même
composante sont dites \emph{isotopes}. Une isotopie entre $f$ et $g$ est une
application $F : X \times I \to Y \times I$ au-dessus de $I = [0,1]$, finie
étale, valant $f$ en $0$ et $g$ en $1$. Comme $Y' \mapsto Y' \times I$ est
une équivalence entre revêtements finis de $Y$ et de $Y \times I$, $F$ se
factorise par un homéomorphisme $\alpha : X \times I \xrightarrow{\sim} Y'
\times I$ suivi de $p \times \mathrm{id}_I$, pour un revêtement $p : Y' \to Y$
unique à isomorphisme unique près\footnote{La page dit le foncteur
« pleinement fidèle » ; il faut aussi qu'il soit essentiellement surjectif, ce
qui vient de ce que $I$ est contractile. C'est ce qu'elle utilise.}. En
d'autres termes, $f = p f_0$ et $g = p g_0$ avec $f_0, g_0 : X \to Y'$ deux
homéomorphismes isotopes, donc $g_0 = f_0 u$ avec $u \in A^{\circ}_X$, et
$g = f u$. Réciproquement $f u$ est isotope à $f$ si $u$ l'est à l'identité.

\textbf{Proposition 1.} \emph{Deux applications finies étales $f, g : X \to Y$
sont isotopes si et seulement si $g = f u$ avec $u \in A^{\circ}_X$.}

La page pose alors la question de l'unicité : toute classe d'isotopie se
factorise trivialement en un homéomorphisme suivi d'un revêtement,
\[
(10)\qquad X \xrightarrow[\ \sim\ ]{\ f_0\ } Y' \xrightarrow{\ p\ } Y ,
\]
mais cette factorisation est-elle canonique, $Y'$ étant défini à
$Y$-isomorphisme unique près ? Il faut pour cela que, pour tout
$Y$-automorphisme $\gamma$ de $Y'$,
\[
(11)\qquad \gamma f_0 \simeq f_0 \ (\text{isotopie}) \Longrightarrow \gamma = \mathrm{id}_{Y'} .
\]
Quitte à prendre $f_0 = \mathrm{id}$, la question est de savoir si un
automorphisme $\gamma$ d'ordre fini, isotope à l'identité, est
l'identité\footnote{La phrase qui pose la question (page 3) se poursuit à la
page 5 ; la page 4 fait suite à la page 5, et la page 6 à la page 4. On lit
dans cet ordre. Le début de la page 4, qui ramène le cas d'une action non libre
au cas d'une action libre en ôtant les points de ramification, est en partie
illisible.}. Ce n'est pas toujours vrai : sur $\mathbb{C}^{*}$, ou sur le tore
$S^1 \times S^1$, les translations par des points d'ordre fini sont isotopes à
l'identité sans être l'identité\footnote{La page écrit
$\mathbb{U} \times \mathbb{U}$, lecture incertaine ; $\mathbb{U}$ est
vraisemblablement le cercle unité, et l'exemple le tore. Les deux exemples sont
de caractéristique d'Euler nulle, hors du cas anabélien.}. Mais :

\textbf{Proposition 2.} \emph{Soit $Y'$ une surface orientée de type fini,
anabélienne. Tout automorphisme d'ordre fini de $Y'$ isotope à l'identité est
l'identité.}

La page indique la voie : on peut supposer $Y'$ connexe ; il suffit que
$\gamma$ agisse trivialement sur $H_1(Y', \mathbb{Z})$ ; on peut munir $Y'$
d'une structure de courbe algébrique invariante par $\gamma$ ; alors $Y'$ se
plonge dans sa jacobienne généralisée — extension d'une variété abélienne par
un tore —, $\gamma$ y induit une translation, et une courbe anabélienne n'est
stable par aucune translation non nulle, « fait bien connu ». Voici un
argument complet. Un homéomorphisme d'ordre fini d'une surface préserve une
structure conforme (Kerékjártó, Eilenberg) ; on peut donc supposer $\gamma$
holomorphe sur $Y' = \overline{Y} \setminus D$, $\overline{Y}$ compacte de
genre $g$, $D$ de cardinal $\nu$, et $\gamma$ s'étend à $\overline{Y}$.
Isotope à l'identité, $\gamma$ agit trivialement sur $H_1(Y', \mathbb{Z})$.
Si $\nu \geq 2$, les classes des petits lacets autour des points de $D$ sont
non nulles et deux à deux distinctes (pour $\nu = 2$, elles sont opposées),
et $\gamma$ les permute : il fixe donc chaque point de $D$. Si $\gamma \neq
\mathrm{id}$, ses points fixes sur $\overline{Y}$ sont isolés et d'indice $+1$,
et la formule de Lefschetz, $\gamma$ agissant trivialement sur
$H_1(\overline{Y})$, donne $2 - 2g = \#\mathrm{Fix}(\gamma) \geq \nu$, soit
$2g - 2 + \nu \leq 0$ ; pour $\nu \leq 1$ elle donne de même $g \leq 1$, avec
$g = 0$ si $\nu = 1$. Dans tous les cas la surface n'est pas anabélienne :
contradiction\footnote{Cet argument est de l'édition. La page dit qu'il suffit
que $\gamma$ agisse trivialement sur $H_1$ à coefficients dans $\mathbb{Z}$,
ou dans $\mathbb{Z}/n\mathbb{Z}$ avec $n \geq 3$ ; ce second énoncé est, pour
les courbes projectives de genre $\geq 2$, le lemme de Serre (1960). La
condition $2g + \nu \geq 2$ que la page écrit (lecture incertaine) est celle
pour que l'application d'Abel–Jacobi généralisée soit injective. Elle ajoute :
« Sûrement, il ne doit pas être difficile de trouver une démonstration directe
purement topologique ».}.

Un revêtement fini d'une surface anabélienne est anabélien. D'où :

\textbf{Corollaire 1.} \emph{Si $Y$ est anabélienne, toute classe d'isotopie
d'applications finies étales $X \to Y$ est un torseur à droite sous
$A^{\circ}_X$.}

En effet l'action $f \mapsto f u$ est transitive par la proposition 1, et
libre : si $f u = f$, $u$ est un automorphisme du revêtement $X \to Y$, donc
d'ordre fini, et isotope à l'identité, donc $u = \mathrm{id}$ par la
proposition 2 appliquée à $X$.

\textbf{Corollaire 2.} \emph{Les composantes connexes de l'espace des
applications finies étales $X \to Y$ sont contractiles.}

C'est le corollaire 1 joint à la contractilité de $A^{\circ}_X$ pour une
surface anabélienne\footnote{La page ne dit pas d'où vient cette
contractilité ; c'est le théorème de Hamstrom (1966) pour les groupes
d'homéomorphismes, d'Earle et Eells (1969) pour les difféomorphismes. Il faut
aussi que l'application orbite $A^{\circ}_X \to$ (composante), $u \mapsto f u$,
soit un homéomorphisme, ce qu'on admet ici avec la page.}.

\textbf{Corollaire 3.} \emph{Le groupe $T_X$ agit à droite sur
$\mathrm{Hom}_{\mathrm{ét}}(X, Y) = \mathrm{Hom}_{\mathrm{isot}}(X, Y)$ avec des
stabilisateurs finis : le stabilisateur de la classe de $f$ s'identifie au
groupe des automorphismes du revêtement $X \to Y$ défini par $f$. Les orbites
sont en bijection avec l'ensemble $R_X(Y)$ des classes d'isomorphisme de
revêtements finis de $Y$ homéomorphes à $X$ ; les orbites de $T^{+}_X$, avec
$R_X(Y) \times \{\pm 1\}$, le signe disant si l'application respecte ou
renverse l'orientation.} L'ensemble des orbites est fini, $\pi_1(Y)$ étant de
type fini et le degré étant fixé par $\chi(X) = d\,\chi(Y)$\footnote{Ces deux
remarques sont dans les marges de la page 7. Une autre note marginale, en
grande partie illisible, rapproche $T_X$ de $\mathrm{Out}$ de $\pi_X$ ; si
c'est bien l'objet de la note, c'est le théorème de Dehn–Nielsen–Baer, mais
la reconstruction est de l'édition.}.

\textbf{Scholie.} Les surfaces orientées de type fini, avec pour flèches les
classes d'isotopie d'applications finies étales qui respectent l'orientation,
forment une catégorie. Une flèche équivaut à un triplet $(Y', p, f_0)$ : un
revêtement fini $p : Y' \to Y$ et une classe d'isotopie d'homéomorphismes
orientés $f_0 : X \to Y'$. À partir de là, tous les morphismes considérés
respectent l'orientation, et l'on omet le signe $+$.

\pagerange{8}{13}
\subsection*{Complétion profinie, et sa version axiomatique (pages 8 à 13)}

On pose
\[
(12)\qquad \widehat{\mathrm{Hom}}_{\mathrm{ét}}(X, Y) = \mathrm{Hom}_{\mathrm{ét}}(X, Y) \times^{T_X} \widehat{T}_X ,
\]
produit contracté par le complété profini $\widehat{T}_X$ de $T_X$. Comme
$T_X$ a un nombre fini d'orbites, à stabilisateurs finis, c'est une réunion
finie d'espaces $\widehat{T}_X/H$, $H$ fini : un espace compact totalement
discontinu, métrisable, sans point isolé\footnote{La page 12 dit d'abord
« compacts », puis « en fait, des espaces localement compacts dénombrables à
l'infini » ; pour $X$, $Y$ fixés, le nombre d'orbites est fini (note de la
page 7) et l'espace est compact. Le « sans point isolé » est son « et pas un
point isolé ».}. La composition s'étend par continuité et reste associative.

\textbf{Version axiomatique.} Une catégorie $C$ est \emph{à groupes
d'automorphismes profinis} si chaque $\mathrm{Aut}(X)$ est muni d'une
structure de groupe profini, de sorte que, pour l'action
$(\psi, \varphi) \cdot u = \varphi u \psi^{-1}$ de $\mathrm{Aut}(X) \times
\mathrm{Aut}(Y)$ sur $\mathrm{Hom}(X, Y)$, les stabilisateurs soient fermés.
On munit alors chaque orbite de la topologie quotient et l'ensemble des
orbites de la topologie discrète : $\mathrm{Hom}(X, Y)$ devient localement
compact. On veut la composition
\[
(14)\qquad \mathrm{Hom}(X, Y) \times \mathrm{Hom}(Y, Z) \longrightarrow \mathrm{Hom}(X, Z)
\]
continue, ce qui revient à la continuité, pour $u$, $v$ fixés, de
$\varphi \mapsto v \varphi u$ sur $\mathrm{Aut}(Y)$. Soit $G_u = \{(\psi,
\varphi) \mid \varphi u = u \psi\}$ le stabilisateur de $u$, et $H_u$ son
image dans $\mathrm{Aut}(Y)$. Sur $H_u$, $v \varphi u = (v u)\psi$ est
continue, la topologie de $H_u$ étant quotient de celle de $G_u$ ; donc aussi
sur chaque translaté $\alpha H_u$. D'où la condition :

\[
(16)\ (*)\qquad \text{pour toute flèche } u : X \to Y,\ H_u \text{ est d'indice fini dans } \mathrm{Aut}(Y).
\]

Fermé d'indice fini, $H_u$ est alors ouvert, ses translatés recouvrent
$\mathrm{Aut}(Y)$, et (14) est continue. La condition dit aussi que, dans
$\mathrm{Hom}(X, Y)$, chaque orbite sous $\mathrm{Aut}(X) \times
\mathrm{Aut}(Y)$ est réunion finie d'orbites ouvertes sous
$\mathrm{Aut}(X)$\footnote{C'est la note marginale de la page 10, en partie
illisible ; l'équivalence est immédiate, l'orbite de $u$ sous
$\mathrm{Aut}(X)$ correspondant à $\mathrm{Aut}(X) \times H_u$.}.

Soit maintenant $C_0$ une catégorie sans topologie, satisfaisant $(*)$ pour
ses groupes d'automorphismes. La page en construit une « complétion
profinie » $\widehat{C}_0$, de mêmes objets, avec
\[
(17)\qquad \mathrm{Hom}_{\widehat{C}_0}(X, Y) = \mathrm{Hom}_{C_0}(X, Y) \times^{\mathrm{Aut}(X) \times \mathrm{Aut}(Y)} \bigl(\widehat{\mathrm{Aut}}(X) \times \widehat{\mathrm{Aut}}(Y)\bigr)
\]
et affirme que ce produit s'identifie à
$\mathrm{Hom}_{C_0}(X, Y) \times^{\mathrm{Aut}(X)} \widehat{\mathrm{Aut}}(X)$,
la composition étant l'unique loi continue qui prolonge celle de
$C_0$\footnote{L'identification n'est pas démontrée par la page. Sur une orbite
$O = (\mathrm{Aut} X \times \mathrm{Aut} Y)/G_u$, elle demande que l'adhérence
de $G_u$ dans $\widehat{\mathrm{Aut}}(X) \times \widehat{\mathrm{Aut}}(Y)$
rencontre $\widehat{\mathrm{Aut}}(X) \times \{1\}$ selon l'adhérence du
stabilisateur de $u$ dans $\mathrm{Aut}(X)$ : une compatibilité des topologies
profinies le long de $u$ que $(*)$ seule ne donne pas. Dans le cas des
surfaces, le stabilisateur est fini, mais l'image de $G_u$ dans $T_X$ — les
classes qui descendent le long du revêtement — est en général d'indice infini,
et la question touche à la topologie que $\widehat{T}_X$ induit sur ce
sous-groupe. On rapporte l'énoncé sans le reprendre à notre compte.}. Tout
$\widehat{C}_0$-morphisme se factorise
\[
(19)\qquad \hat f : X \xrightarrow[\ \sim\ ]{\ \hat f_0\ } Y' \xrightarrow{\ u\ } Y ,
\]
$u$ morphisme de $C_0$, $\hat f_0$ isomorphisme de $\widehat{C}_0$ ; si
$\mathrm{Aut}_Y(Y') \to \widehat{\mathrm{Aut}}(Y')$ est injectif pour toute
$u : Y' \to Y$, $Y'$ est déterminé à $Y$-isomorphisme unique près. Pour
composer $\hat f = u \hat f_0$ et $\hat g = v \hat g_0$, $\hat g_0 : Y
\xrightarrow{\sim} Z'$, on écrit $\hat g_0 = w \hat\varphi$ avec $w$ un
isomorphisme de $C_0$ et $\hat\varphi \in \widehat{H}_u$ — c'est possible
parce que $\widehat{H}_u$ est ouvert —, on choisit $\hat\varphi'$ avec
$\hat\varphi u = u \hat\varphi'$, et
\[
(22)\qquad \hat g \hat f = (v w \hat\varphi)(u \hat f_0) = \underbrace{(v w u)}_{\text{flèche de } C_0}\ \underbrace{(\hat\varphi' \hat f_0)}_{\text{isomorphisme de } \widehat{C}_0} .
\]
Dans le cas des surfaces anabéliennes, $C_0$ est la catégorie isotopique, la
condition $(*)$ est satisfaite, et $\widehat{C}_0$ est la catégorie définie par
(12)\footnote{La page 13 fait suite à la page 11 (« Par construction, tout
$\widehat{C}_0$-morphisme » — « (18) $\hat f : X \to Y$ se factorise »), et la
page 14 à la page 12. On lit dans cet ordre ; la transcription suit l'ordre des
feuillets.}.

\pagerange{12}{16}
\subsection*{La catégorie de Teichmüller–Galois (pages 12, 14 à 16)}

Limitons-nous aux surfaces qui sont des courbes algébriques complexes
anabéliennes. Dans la factorisation (19), le revêtement $u : Y' \to Y$ provient
d'un revêtement fini étale \emph{algébrique}, et l'isomorphisme $\hat f_0$
s'interprète comme une flèche du groupoïde fondamental profini du champ des
modules pertinent — la « multiplicité modulaire » $M_{g,\nu}$ — entre les
points définis par $X$ et $Y'$ : une classe d'isotopie d'homéomorphismes
$X \to Y'$ est exactement une classe de chemins de $[X]$ à $[Y']$ dans ce
champ, dont l'espace de Teichmüller est le revêtement universel. La
composition (22) s'explicite donc en compositions dans le groupoïde
fondamental du champ des modules et dans la catégorie des courbes algébriques
avec leurs revêtements. Tout devient algébrique, et la construction a un sens
sur un corps algébriquement clos $k$ quelconque ; elle est fonctorielle en $k$,
les foncteurs de changement de base $k \to k'$ étant des équivalences, parce
que le groupoïde fondamental du champ des modules ou d'une courbe ne change pas
par extension de corps algébriquement clos\footnote{La page ajoute « pas besoin
qu'il soit de car.~0 ! ». En caractéristique $p > 0$, l'invariance par
extension de corps algébriquement clos vaut pour les variétés propres et pour
les quotients modérés, mais pas pour le groupe fondamental d'une courbe
affine, que la ramification sauvage fait dépendre du corps ; et les
revêtements étales d'une telle courbe ne relèvent plus de la topologie. On
s'en tient ici à la caractéristique nulle. La note marginale de la page 15
précise que l'équivalence ne signifie pas que toute courbe sur $k'$ provienne
de $k$ : c'est que, dans cette catégorie, deux courbes de même type sont
isomorphes.}.

En particulier $\mathrm{Aut}(k)$ agit sur cette catégorie, qu'il nomme
\emph{catégorie de Teichmüller–Galois} associée à $k$. Ses objets connexes à
isomorphisme près correspondent aux types $(g, \nu)$ anabéliens ; le groupe
des automorphismes d'un objet est le complété profini du groupe de
Teichmüller, c'est-à-dire le groupe fondamental profini du champ
$M_{g,[\nu]}$ des courbes de genre $g$ à $\nu$ points non ordonnés\footnote{La
fin de la page 15 est peu lisible ; « des compactifiés profinis des groupes
… de Teichmüller » est sûr, le reste non.}. Élargir le \emph{groupoïde} de
Teichmüller en cette \emph{catégorie}, qui n'est plus un groupoïde, est une
façon de relier entre eux les types $(g, \nu)$ différents.

« Une part essentielle du vaste programme à l'horizon est la détermination
complète de cette catégorie par voie “algébrique” (ni transcendante, ni de
géométrie algébrique !), avec les opérations de Galois–Teichmüller dessus »,
au moins pour $k = \overline{\mathbb{Q}}$, et dans une certaine mesure pour
$k = \overline{K}$, $K$ de type fini sur $\mathbb{Q}$, en se limitant aux
$K$-automorphismes. « Algébrique » veut dire ici, précise la marge, par des
techniques combinatoires sur les surfaces et les groupes
profinis\footnote{Note marginale très serrée, lecture en grande partie
conjecturale. Ce programme est celui que l'\emph{Esquisse d'un programme}
(1984) exposera sous le nom de « tour de Teichmüller ».}.

\pagerange{17}{20}
\subsection*{Le bouchage des trous (pages 17 à 20)}

Un type de morphismes manque : les inclusions ouvertes
\[
(23)\qquad X = Y \setminus T \hookrightarrow Y ,
\]
qui \emph{bouchent des trous}. Pour les composer avec les morphismes finis
étales, il faudrait admettre des morphismes $X \to Y$ qui s'insèrent dans un
carré
\[
(24)\qquad
\begin{array}{ccc}
X & \hookrightarrow & \widehat{X}\\
\downarrow & & \downarrow\\
Y & \hookrightarrow & \widehat{Y}
\end{array}
\]
avec $\widehat{X} \to \widehat{Y}$ fini ramifié et $X \to Y$ étale. Si $S =
\widehat{X} \setminus X$ et $T = \widehat{Y} \setminus Y$, on a
$S \supset \widehat{X}|_T$, l'image inverse de $T$\footnote{La page écrit
$T = \widehat{Y} \setminus T$, pour $\widehat{Y} \setminus Y$.}, et la
difficulté vient de ce que $\widehat{X} \to \widehat{Y}$ peut être ramifié non
seulement en des points de $\widehat{X}|_T$, mais en des points de $S$ qui
n'y sont pas — des trous de $X$ au-dessus de points de $Y$. La notion
d'isotopie raisonnable existe (dans $\mathrm{Aut}(\widehat{X}, S) \times
\mathrm{Aut}(\widehat{Y}, T)$), mais il n'y a plus de factorisation
\emph{canonique} d'une classe en un bouchage de trous suivi d'un revêtement :
les points de $S \setminus \widehat{X}|_T$ vont en des points de $Y$ que la
classe seule ne détermine pas. « Pour cette raison, je ne vois plus comment
interpréter une telle classe […] en termes de géométrie algébrique
“abstraite” »\footnote{La page 18 est très corrigée ; l'ordre de lecture de
sa phrase centrale est conjectural, et l'on n'en retient que la conclusion,
qui est sûre.}.

On peut pourtant définir, sur tout corps algébriquement clos, un foncteur du
groupoïde de Teichmüller des courbes de type $(g, \nu)$ munies d'un ensemble
$I$ de $\nu' \leq \nu$ trous marqués vers celui du type $(g, \nu - \nu')$, par
« bouchage des trous en $I$ », transitif pour des bouchages
successifs\footnote{Aujourd'hui : l'application d'oubli de points
$M_{g,\nu} \to M_{g,\nu-\nu'}$ des champs de modules, sur les groupoïdes
fondamentaux. Pour rester dans le monde anabélien, il faut $2g - 2 + \nu -
\nu' > 0$, ce que la page ne précise pas.}. Ce qu'il ne voit pas, c'est
comment loger ces inclusions dans une même catégorie que les morphismes
étales — sinon en définissant par générateurs et relations une catégorie plus
grosse, au prix de ne plus pouvoir l'interpréter, sur $\mathbb{C}$, comme
complétée profinie de la catégorie isotopique. Le développement s'arrête là,
au milieu de la page 20. La suite V du dossier (pages 120 et 121) traitera le
bouchage des trous sur les présentations combinatoires, où il devient trivial.

\pagerange{21}{41}
\section*{II. Opérations de $\Gamma_{\overline{\mathbb{Q}}/\mathbb{Q}}$ sur $\pi_1 M_{0,3}$ (pages 21 à 41)}

La page 21 est une couverture qui porte ce titre ; la page 22 n'a qu'un
décompte au crayon. Le texte commence à sa page 5, les formules repartant de
(1) : ses pages 1 à 4, qui fixaient les conventions, manquent. Les calculs de
cette suite sont donc rapportés tels que les pages les donnent ; seules les
implications entre équations qui ne dépendent pas de ces conventions ont été
vérifiées, et on le signale\footnote{« $\pi_1 M_{0,3}$ » désigne ici le groupe
fondamental d'une courbe de type $(0, 3)$, c'est-à-dire $\Pi_{0,3}$, et non
celui du champ des modules $M_{0,3}$, qui est un point ; c'est l'usage de la
\emph{Longue Marche}, qui écrit $\pi_{0,3}$.}.

\emph{Le cadre.} Sur la sphère, les trois points marqués $R_0, R_1, R_\infty$
sont sur l'équateur ; $\underline{\omega}$ est l'ensemble à deux éléments des
ordres cycliques de $J = \{0, 1, \infty\}$ (les deux sens de parcours de
l'équateur), et pour $\omega \in \underline{\omega}$ on note $\omega i$
l'image de $i$ par la permutation circulaire correspondante, $\omega'$ l'autre
élément. Le groupe $\mathfrak{S}_J$ agit sur la sphère par rotations — le
tiers de tour autour de l'axe des pôles, les demi-tours autour des axes qui
passent par un $R_i$ —, et $\tau$ est la réflexion par rapport au plan de
l'équateur\footnote{La page note ce groupe $\widetilde{\mathfrak{S}}_J$ et dit
qu'il permute « de façon simplement transitive » chacun des paquets de six
flèches introduits ensuite, y compris les six méridiens $d_i^{\omega}$ issus
des deux pôles. Cela impose l'action par rotations, les demi-tours échangeant
les pôles ; un groupe fixant un pôle ne pourrait envoyer $d_i^{+}$ sur
$d_i^{-}$. L'identification est de l'édition.}. On cherche les automorphismes
$u$ du complété profini de divers groupoïdes fondamentaux de la sphère
trouée, qui \emph{fixent les sommets} et \emph{commutent à $\mathfrak{S}_J$} :
ce sont les candidats à l'action d'un élément de Galois, qui fixe les points
réels et commute aux symétries définies sur $\mathbb{Q}$.

\pagerange{23}{24}
\subsection*{Le groupoïde $\mathrm{I}_J$ (pages 23 et 24)}

Autour de chaque trou $R_i$, un petit cercle porte deux sommets
$R_i^{\omega}$, $\omega \in \underline{\omega}$, reliés par deux arcs
$a_i^{\omega} : R_i^{\omega} \to R_i^{\omega'}$ ; le long de l'équateur, des
flèches $b_i^{\omega} : R^{\omega}_{\omega i} \to R^{\omega'}_{\omega' i}$
relient les sommets voisins de deux trous différents. Le groupoïde
$\mathrm{I}_J$ a ces six sommets, ces douze flèches, et les relations
\[
(2)\qquad \text{a)}\ b_i^{\omega} b_i^{\omega'} = 1, \qquad
\text{b)}\ \text{une relation d'hexagone par hémisphère},
\]
soit trois relations effectives en a) et deux en b)\footnote{Les indices de
(2\,b) sont lus avec doute ; on n'en retient que la forme. Le compte est
cohérent : $12 - 3 = 9$ flèches indépendantes, $6$ sommets, $2$ relations, et
$9 - 6 + 1 - 2 = 2$ est bien le rang du groupe libre $\Pi_{0,3}$
(vérification de l'édition).}. Le lacet complet $\lambda_i^{\omega} =
a_i^{\omega'} a_i^{\omega}$ autour de $R_i$ est l'image de $l_0$ par un
isomorphisme $\varphi_i^{\omega} : \Pi_{0,3} \xrightarrow{\sim}
\pi(R_i^{\omega})$, les six $\varphi_i^{\omega}$ étant permutés simplement
transitivement par $\mathfrak{S}_J$ ; $\tau$ agit par $\tau_\infty$ (5). Les
arcs transportent ces repères selon
\[
(6)\qquad a_i^{\omega}\bigl(\varphi_i^{\omega}(g)\bigr) = \varphi_i^{\omega'}\bigl(\varepsilon_0(g)\bigr), \qquad
b_i^{\omega}\bigl(\varphi^{\omega}_{\omega i}(g)\bigr) = \varphi^{\omega'}_{\omega' i}\bigl(\sigma_\infty(g)\bigr), \qquad \varepsilon_0^2 = l_0 ,
\]
où $\varepsilon_0$ est la « racine carrée » de $l_0$ de la page 54 : faire le
demi-tour d'un trou, c'est appliquer $\varepsilon_0$ ; le tour complet, deux
demi-tours, c'est conjuguer par $l_0$\footnote{Avec $\varepsilon_0 =
\sigma_\infty \rho$ (page 54), on vérifie $\varepsilon_0(l_0) = l_0$ et
$\varepsilon_0^2 = \mathrm{int}(l_0)$ comme automorphismes de $\Pi_{0,3}$ ;
c'est ce qu'utilise la formule (30) ci-dessous. Le dossier 145 écrit de même
$\varepsilon_i^2 = l_i$, ses $\varepsilon_i$ relevant les transpositions.}.
Avec $\tilde b_i^{\omega} = a^{\omega'}_{\omega' i} b_i^{\omega}$, la relation
(2\,b) devient
\[
(8)\qquad \tilde b^{\omega}_{\omega^2 i}\, \tilde b^{\omega}_{\omega i}\, \tilde b_i^{\omega} = 1 ,
\]
et $\tilde b$ transporte par $\rho^{-1}$ (9).

Un automorphisme $u$ du complété $\widehat{\mathrm{I}}_J$ commutant à
$\mathfrak{S}_J$, fixant les sommets et respectant les sous-groupoïdes
engendrés par chaque petit cercle, s'écrit
\[
u(a_i^{\omega}) = a_i^{\omega}\, \varphi_i^{\omega}(l_0^{\gamma}), \qquad
u(b_i^{\omega}) = b_i^{\omega}\, \varphi^{\omega}_{\omega i}(\beta),
\qquad \gamma = \frac{p - 1}{2} \in \widehat{\mathbb{Z}},\ p \in \widehat{\mathbb{Z}}^{*},\ \beta \in \widehat\Pi_{0,3} ;
\]
alors $u(\lambda_i^{\omega}) = (\lambda_i^{\omega})^{p}$ : $p$ est le
\emph{multiplicateur}, l'exposant par lequel $u$ élève les lacets autour des
trous. Les relations (2) donnent, sur $(p, \beta)$,
\[
(12)\qquad \beta\, \sigma_\infty(\beta) = 1, \qquad
\alpha\, \rho^{-1}(\alpha)\, \rho^{-2}(\alpha) = 1 \quad \text{où } \alpha = l_1^{\gamma}\, \beta = \sigma_\infty(l_0^{\gamma})\,\beta .
\]
La marge remarque qu'en passant aux inverses on retrouve « les équations
habituelles » $\alpha'\rho(\alpha')\rho^{2}(\alpha') = 1$, $\alpha' =
\alpha^{-1}$ (c'est exact : on applique $\rho^{2}$ à l'inverse de la
relation). Le passage de $u$ à $\tau u \tau^{-1}$ change $\beta$ en
$\tau_\infty(\beta)$, et il faudrait, dit-il, vérifier directement que (12)
est stable par ce changement.

\emph{Le nom d'aujourd'hui.} Ces deux équations ont exactement la forme des
deux premières relations par lesquelles Drinfeld (1990) définit le
\emph{groupe de Grothendieck–Teichmüller} $\widehat{GT}$ : pour
$(\lambda, f) \in \widehat{\mathbb{Z}}^{*} \times \widehat{F}_2$,
$f(x, y) f(y, x) = 1$ et $f(z, x) z^{m} f(y, z) y^{m} f(x, y) x^{m} = 1$ avec
$xyz = 1$ et $m = (\lambda - 1)/2$. Ici $\sigma_\infty$ échange $l_0$ et $l_1$,
$\rho$ les permute circulairement, et le multiplicateur $p$ entre par
$\gamma = (p - 1)/2$ comme $\lambda$ par $m$\footnote{Le rapprochement est de
l'édition, à des conventions près (ordre de la relation entre les lacets,
côté des conjugaisons) que nous n'avons pas confrontées terme à terme. Ni
Drinfeld ni Ihara, qui nomma le groupe (1991), ne pouvaient être cités par ces
pages. La troisième relation de Drinfeld, pentagonale, porte sur $M_{0,5}$ ;
voir la page 99 ci-dessous. Le dossier 145 rencontre ces mêmes relations, dans
une autre présentation, aux pages 6 à 11.}.

\pagerange{25}{28}
\subsection*{Les groupoïdes $\mathrm{II}_J$ et $\mathrm{III}_J$ (pages 25 à 28)}

On ajoute sur l'équateur les trois sommets $Q_i$, opposés aux trous, et des
arcs $c_i^{\omega} : Q_i \to R^{\omega'}_{\omega' i}$. Alors $b_i^{\omega} =
c_i^{\omega}(c_i^{\omega'})^{-1}$ (13), les $b$ deviennent superflus, (2\,a)
est automatique, et avec
\[
(14)\qquad \tilde a_i^{\omega} = (c^{\omega}_{\omega i})^{-1}\, a_i^{\omega}\, c^{\omega'}_{\omega' i} : Q_{\omega i} \to Q_{\omega' i}, \qquad \tau(\tilde a_i^{\omega}) = (\tilde a_i^{\omega'})^{-1},
\]
(2\,b) devient $\tilde a^{\omega}_{\omega^2 i}\, \tilde a^{\omega}_{\omega i}\,
\tilde a_i^{\omega} = 1$ (16). On repère $\pi(Q_i)$ par $\psi_i^{\omega} :
\Pi_{0,3} \xrightarrow{\sim} \pi(Q_i)$, déduit de
$\varphi^{\omega'}_{\omega' i}$ par transport le long de $c_i^{\omega}$.

Un automorphisme $u$ de $\widehat{\mathrm{II}}_J$ du type considéré est donné
par $u(a_i^{\omega}) = a_i^{\omega}\varphi_i^{\omega}(\lambda)$ et
$u(c_i^{\omega}) = c_i^{\omega}\psi_i^{\omega}(\eta)$, $\lambda, \eta \in
\widehat\Pi_{0,3}$ — la marge observe qu'il est inutile d'exiger d'emblée
$\lambda = l_0^{\gamma}$. Alors $u(b_i^{\omega}) =
b_i^{\omega}\varphi^{\omega}_{\omega i}(\beta)$ avec
\[
(23)\qquad \beta = \sigma_\infty(\eta)\, \eta^{-1},
\]
qui entraîne la première relation (12), $\beta\sigma_\infty(\beta) = 1$
($\sigma_\infty$ étant involutif) ; il reste la seconde, (24). Sur les
$\tilde a$, $u$ agit par
\[
(27)\qquad \tilde\alpha = \eta^{-1}\, \varepsilon_0(\lambda\eta), \qquad \text{soit } \eta^{-1} l_0^{\gamma}\, \varepsilon_0(\eta) \text{ si } \lambda = l_0^{\gamma},
\]
et (24) équivaut à la même équation en $\tilde\alpha$ :
\[
(26)\qquad \tilde\alpha\, \rho^{-1}(\tilde\alpha)\, \rho^{-2}(\tilde\alpha) = 1 .
\]
Le groupoïde $\mathrm{III}_J$ induit sur les trois $Q_i$ (six arcs
$\tilde a_i^{\omega}$, relations (16)) a donc pour automorphismes du type
considéré les $\tilde\alpha \in \widehat\Pi_{0,3}$ solutions de (26). Sur le
lacet $\tilde\lambda_i^{\omega} = \tilde a_i^{\omega'}\tilde a_i^{\omega}$,
identifié à $l_0$,
\[
(30)\qquad u(l_0) = \tilde\alpha\, \varepsilon_0(\tilde\alpha)\, l_0 = \eta^{-1}\, l_0^{p}\, \eta \qquad (p = 2\gamma + 1) :
\]
$u$ envoie le lacet autour d'un trou sur un conjugué de sa puissance $p$-ième
— c'est la propriété qui définit les automorphismes « à lacets » du dossier
145, ceux qu'on dit aujourd'hui \emph{préserver l'inertie}\footnote{La
seconde égalité de (30) a été vérifiée, à partir de $\varepsilon_0(l_0) = l_0$
et $\varepsilon_0^2 = \mathrm{int}(l_0)$.}.

Le passage de $u$ à $\tau u \tau^{-1}$ change $\tilde\alpha$ en
$\varepsilon_0\tau_\infty(\tilde\alpha^{-1})$. La page en tire une remarque
générale, juste : le groupe $\widetilde{\mathfrak{S}}_3$ engendré par $\rho$
et les $\tilde\sigma_i$, qui contient $\{1, \rho, \rho^2\}$ avec indice $2$,
agit sur l'ensemble des solutions $x$ de $x\,\rho^{-1}(x)\,\rho^{-2}(x) = 1$
par
\[
(32)\qquad \Theta_g(x) = g\bigl(x^{\chi(g)}\bigr), \qquad \chi(g) = \pm 1 \text{ la signature},
\]
et $\tau$ agit par $\Theta_{\tilde\sigma_0}$ (33)\footnote{(32) est vérifié
ici : pour $g = \rho$, on applique $\rho$ à la relation ; pour $g$ impair, qui
vérifie $g\rho = \rho^{-1}g$, l'expression transformée est l'image par $g$ de
l'inverse de $\rho^{-1}(x)\rho^{-2}(x)x$, conjugué cyclique de
$x\rho^{-1}(x)\rho^{-2}(x)$, donc trivial. L'exposant $\chi(g)$ est surchargé
sur la page.}.

\pagerange{29}{31}
\subsection*{Les groupoïdes $\mathrm{IV}_J$ et $\mathrm{V}_J$ : les pôles (pages 29 à 31)}

On ajoute les deux pôles $P^{\omega}$ et les six méridiens $d_i^{\omega} :
P^{\omega} \to Q_i$, avec $\tau(P^{\omega}) = P^{\omega'}$,
$\tau(d_i^{\omega}) = d_i^{\omega'}$. Le groupoïde $\mathrm{IV}_J$ (sommets
$Q_i$, $P^{\omega}$ ; arcs $\tilde a$, $d$ ; relations (35)) est équivalent à
$\mathrm{III}_J$, et ses automorphismes du type considéré sont donnés par
$\tilde\alpha$ et $\zeta \in \widehat\Pi_{0,3}$, liés par
\[
(37)\qquad \tilde\alpha = \zeta\, \rho^{-1}(\zeta^{-1}) .
\]
\emph{Il n'y a plus de condition sur $\zeta$} : toute $\zeta$ convient, et
(37) entraîne (26) parce que $\rho$ est d'ordre trois\footnote{Vérifié : le produit
$\tilde\alpha\rho^{-1}(\tilde\alpha)\rho^{-2}(\tilde\alpha)$ se télescope en
$\zeta\rho^{-3}(\zeta)^{-1} = 1$.}. C'est le gain des pôles : ce sont les
points fixes de la rotation $\rho$, et pris pour points base ils
« résolvent » l'équation d'ordre trois. La marge note que les arcs $\tilde
a_i^{\omega}$ deviennent superflus.

Le groupoïde $\mathrm{V}_J$ induit sur les deux pôles est libre, engendré par
trois arêtes $\tilde d_i^{\omega} = (d_i^{\omega'})^{-1} d_i^{\omega} :
P^{\omega} \to P^{\omega'}$ ; ses automorphismes sont donnés par un
$\tilde\beta \in \widehat\Pi_{0,3}$ soumis à la seule relation
\[
(47)\qquad \tilde\beta\, \sigma_\infty(\tilde\beta) = 1 ,
\]
et $\tilde\beta = \zeta\sigma_\infty(\zeta)^{-1}$ quand $u$ vient de
$\mathrm{IV}_J$.

\pagerange{31}{36}
\subsection*{Le groupoïde $\mathrm{VI}$ et la relation des lacets (pages 31 à 36)}

Le groupoïde induit sur les onze sommets $R_i^{\omega}$, $Q_i$, $P^{\omega}$
est engendré par les arcs $a$, $c$, $d$ — trois paquets homogènes sous
$\mathfrak{S}_J$ — avec six relations (48). Un automorphisme est donné par
\[
(54)\qquad u(a_i^{\omega}) = a_i^{\omega}\varphi_i^{\omega}(\lambda), \qquad
u(c_i^{\omega}) = c_i^{\omega}\psi_i^{\omega}(\eta), \qquad
u(d_i^{\omega}) = d_i^{\omega}\tilde\psi_i^{\omega}(\zeta),
\]
et, au lieu de refaire le calcul, il observe que $\mathrm{VI}$ est la réunion
de $\mathrm{II}$ et de $\mathrm{IV}$, d'intersection $\mathrm{III}$ : un
automorphisme de $\mathrm{VI}$ est un couple d'automorphismes de $\mathrm{II}$
et de $\mathrm{IV}$ qui se recollent sur $\mathrm{III}$, c'est-à-dire tels que
les deux expressions de $\tilde\alpha$ coïncident :
\[
(56)\qquad \eta^{-1}\varepsilon_0(\lambda\eta) = \zeta\, \rho^{-1}(\zeta)^{-1} .
\]
C'est la seule relation entre $\lambda, \eta, \zeta$ : celle de
$\mathrm{II}$ est conséquence de celle de $\mathrm{IV}$. Pour la mettre
« sous une forme plus familière pour moi », il passe aux inverses,
$\beta' = \beta^{-1}$, $\alpha' = \alpha^{-1}$, $\tilde\alpha' =
\tilde\alpha^{-1}$, et pose $\zeta' = \eta\zeta$. En résumé :

\emph{Les automorphismes de $\widehat{\mathrm{VI}}_J$ qui fixent les sommets et
commutent à $\mathfrak{S}_J$ correspondent aux triplets $(\lambda, \eta,
\zeta') \in \widehat\Pi_{0,3}^{\,3}$ satisfaisant la \emph{relation des
lacets}}
\[
(63)\qquad \boxed{\ \beta' = \alpha'\, \sigma_\infty(\lambda), \qquad \alpha' = \zeta'\, \rho(\zeta')^{-1}, \qquad \beta' = \eta\, \sigma_\infty(\eta)^{-1}\ }
\]
avec $\sigma_\infty(\lambda) = l_1^{\gamma}$ quand $\lambda =
l_0^{\gamma}$, dans la forme qu'on retient ici\footnote{La page écrit, en (60) et (63),
$\alpha' = \zeta'\rho^{-1}(\zeta')^{-1}$, l'exposant de $\rho$ étant surchargé
et lu $-1$. Mais la dérivation qu'elle donne à la page 35 — (25\,bis) passé
aux inverses, $\alpha' = \eta\,\rho(\tilde\alpha')\,\rho(\eta)^{-1}$, puis
$\rho(\tilde\alpha') = \zeta\rho(\zeta)^{-1}$ — aboutit à
$\alpha' = (\eta\zeta)\,\rho(\eta\zeta)^{-1}$, et c'est ce que réécrit (78) à
la page 40. C'est aussi la forme qui donne la relation « habituelle »
$\alpha'\rho(\alpha')\rho^2(\alpha') = 1$ par télescopage. On l'adopte. Les
autres pas de (57) à (63) ont été vérifiés ; la formule (25\,bis), qui dépend
des conventions perdues, est reprise de la page.}.

Les quantités $\alpha', \beta', \zeta'$ ont été introduites « par le pur
calcul », comme donnant les formules les plus sympathiques ; il leur donne
ensuite un sens géométrique, en renversant les flèches $b$, $\tilde b$,
$\tilde a$ (64)–(65) et en introduisant les chemins
$\tilde d_i'^{\omega} = c^{\omega}_{\omega i}\, d^{\omega}_{\omega i} :
P^{\omega} \to R_i^{\omega'}$, du pôle au voisinage du trou, sur lesquels
\[
(67)\qquad u(\tilde d_i'^{\omega}) = \varphi_i^{\omega'}(\zeta')\, \tilde d_i'^{\omega} :
\]
$\zeta'$ mesure ce que $u$ fait au chemin direct d'un pôle à un trou.

\pagerange{37}{41}
\subsection*{Le groupoïde $\mathrm{VII}$, et le diagramme des sept groupoïdes (pages 37 à 41)}

Deux cas restent : le groupoïde $\mathrm{VII}$ induit sur les huit sommets
$R_i^{\omega}$, $P^{\omega}$, et celui induit sur les $Q_i$, $P^{\omega}$, qui
n'est autre que $\mathrm{IV}$ présenté par les seuls méridiens (les
$\tilde a$ y sont inutiles par (35)). Les sept groupoïdes s'emboîtent selon le
diagramme (68)\footnote{La page encadre ensemble $\mathrm{I}$, $\mathrm{II}$,
$\mathrm{VII}$, $\mathrm{VI}$ ; une flèche partant de $\mathrm{II}$ vers la
droite est rayée. Le diagramme a sept nœuds et huit inclusions, tous
reproduits.} :

\begin{tikzcd}[column sep=small, row sep=small]
& & \mathrm{III} \arrow[dl, hook'] \arrow[dr, hook] & & \\
\mathrm{I} \arrow[r, hook] \arrow[dr, hook] & \mathrm{II} \arrow[dr, hook] & & \mathrm{IV} \arrow[dl, hook'] & \mathrm{V} \arrow[l, hook'] \\
& \mathrm{VII} \arrow[r, hook] & \mathrm{VI} & &
\end{tikzcd}

Pour $\mathrm{VII}$, le calcul prédit que ses automorphismes correspondent aux
couples $(\lambda, \zeta')$ — $\beta'$ s'en déduit, d'où la marge « donc
$\beta'$ inutile ! » — soumis à
\[
(70)\qquad \beta' = \alpha'\sigma_\infty(\lambda),\quad \alpha' = \zeta'\rho(\zeta')^{-1}, \qquad \beta'\sigma_\infty(\beta') = 1, \ \text{soit}\ \alpha'\,\sigma_\infty(\lambda)\,\sigma_\infty(\alpha')\,\lambda = 1 ,
\]
tandis que ceux de $\mathrm{II}$ correspondent aux $(\lambda, \eta)$ — « donc
$\alpha'$ inutile ! » — avec $\beta' = \eta\sigma_\infty(\eta)^{-1} =
\alpha'\sigma_\infty(\lambda)$ et $\alpha'\rho(\alpha')\rho^2(\alpha') = 1$
(70$'$)\footnote{Les deux formes de la seconde ligne de (70) sont équivalentes
($\sigma_\infty$ est involutif) ; vérifié. Même remarque qu'en (63) sur
l'exposant de $\rho$.}. Les pages 39 à 41 explicitent le cas de
$\mathrm{VII}$, présenté par les chemins $\tilde d_i^{\omega} : P^{\omega} \to
R_i^{\omega}$, $\tilde d_i'^{\omega} : P^{\omega} \to R_i^{\omega'}$ et les
$a_i^{\omega}$, avec deux familles de relations (72), dont trois seulement
sont indépendantes. On y pose $u(\tilde d) = \varphi(\zeta_1)\tilde d$,
$u(\tilde d') = \varphi(\zeta')\tilde d'$, $u(a) = a\varphi(\lambda)$ ; la
première relation donne $\zeta' = \varepsilon_0(\lambda\zeta_1)$ (74), la
seconde équivaut à $\beta'\sigma_\infty(\beta') = 1$ avec
$\beta' = \alpha'\sigma_\infty(\lambda)$, $\alpha' = \zeta'\rho(\zeta')^{-1}$
(78), et la page 41 conclut sur l'égalité des diverses expressions de
$\beta'$, « ce qui donne la formulation la plus simple »\footnote{Les indices
de (72) sont surchargés ; le numéro (75) sert deux fois, (78) trois fois ; les
formules (76) et (79) auxquelles renvoie la page 41 ne figurent pas au
dossier. On n'en retient que la conclusion, qui coïncide avec (70).}.

\pagerange{43}{63}
\section*{III. La sphère à trois points marqués (pages 43 à 63)}

Une nouvelle suite, paginée de 1 à 12, reprend la même sphère, cette fois sans
rien choisir : chaque objet est construit de façon que les symétries le
transportent sur un objet du même type. Il l'annonce comme « exercice
préliminaire de syntaxe galoisienne », avant des configurations plus riches en
symétries sur la sphère ou sur d'autres surfaces, « en vue des applications au
calcul des groupoïdes de Teichmüller ».

\pagerange{43}{45}
\subsection*{La géométrie (pages 43 à 45)}

Soit $\Sigma$ une sphère conforme avec une partie $J = \{R_i\}$ à trois
éléments. Son groupe d'automorphismes conformes ou anticonformes respectant $J$
est
\[
\Gamma_{\Sigma, J} \simeq \mathfrak{S}_J \times \{1, \tau\},
\]
$\tau$ l'unique antiautomorphisme fixant les $R_i$ — la conjugaison complexe
quand $J = \{0, 1, \infty\}$\footnote{La marge note que l'antipodie n'est pas
dans $\Gamma_{\Sigma,J}$ : elle ne stabilise pas $J$.}. Une métrique
riemannienne invariante, unique à scalaire près, identifie $\Sigma$ à la sphère
unité d'un espace euclidien $E$ de dimension $3$. Le plan $V$ engendré par les
$R_j$ coupe $\Sigma$ selon un grand cercle $\Sigma_{\mathbb{R}}$, sur lequel les
$R_j$ forment un triangle équilatéral : $\sum R_j = 0$, et
\[
(5),(7)\qquad V \simeq V_J(\mathbb{R}) = \ker\bigl(\mathbb{R}^J \xrightarrow{\ \Sigma\ } \mathbb{R}\bigr), \qquad
E \simeq V_J(\mathbb{R}) \oplus \Delta ,
\]
la forme quadratique sur $V_J(\mathbb{R})$ étant celle qui rend les $R_j$
unitaires — $x^2 - xy + y^2$ dans la base $(R_i, R_j)$, où
$R_k = (-1, -1)$ — et $\Delta = V^{\perp}$ l'axe polaire. Trois ensembles à
deux éléments apparaissent :
\[
(3)\qquad \boldsymbol{\varpi} = \text{orientations de } \Sigma \text{ (ou de } E), \qquad
\boldsymbol{\omega} = \text{orientations de } \Sigma_{\mathbb{R}} \text{ (ordres cycliques de } J),
\]
\[
\boldsymbol{\varepsilon} = \pi_0(\Sigma \setminus \Sigma_{\mathbb{R}}) \text{ (hémisphères)},
\]
liés par $\boldsymbol{\varpi} \simeq \boldsymbol{\omega} \wedge
\boldsymbol{\varepsilon}$ : une orientation de l'équateur et le choix d'un
hémisphère orientent la sphère. $\Delta$ est la droite réelle tordue par
$\boldsymbol{\varepsilon}$, ses deux demi-droites étant indexées par les
hémisphères. $\mathfrak{S}_J$ agit sur $E$ par l'action tautologique sur
$V_J(\mathbb{R})$ et par la signature sur $\Delta$ (8) — une transposition, qui
renverse le plan, doit renverser l'axe pour rester une rotation —, et
$\tau = \mathrm{id} \oplus (-\mathrm{id}_\Delta)$ (9)\footnote{Vérifications de
l'édition : $q(R_k) = 1 - 1 + 1 = 1$, et le produit scalaire $-1/2$ entre deux
$R_j$ est celui d'un triangle équilatéral inscrit.}.

\pagerange{45}{48}
\subsection*{L'éclatement réel, et la représentation fictive (pages 45 à 48)}

On s'intéresse au groupoïde fondamental de $\Sigma^{*} = \Sigma \setminus J$.
Pour disposer de points base « proches » de chaque $R_i$, on compactifie : soit
$\Sigma'$ l'éclaté réel de $\Sigma$ en $J$, où chaque $R_i$ est remplacé par la
droite projective réelle $P(T_{\Sigma, R_i})$ de ses directions tangentes. C'est
la sphère où trois disques sont remplacés par des rubans de Möbius, une surface
non orientable de caractéristique d'Euler $2 - 3 = -1$, à trois bonnets croisés\footnote{La page conclut
« donc le genre est $2$ », par $\chi = 1 - g$. Le nombre de bonnets croisés de
cette surface est $3$ ($\chi = 2 - 3$) ; $2$ est son premier nombre de Betti
rationnel. La surface elle-même ne sert qu'à passer à $\widetilde\Sigma$.}. En
la découpant le long des trois cercles, on obtient une surface compacte à bord
$\widetilde\Sigma$ — une sphère à trois trous —, dont le bord
\[
(12)\qquad \partial\widetilde\Sigma \simeq \coprod_{i \in J} U_i , \qquad U_i = \text{cercle des directions tangentes unitaires en } R_i ,
\]
revêt doublement les droites projectives : c'est l'\emph{éclatement réel
orienté}. L'intérieur de $\widetilde\Sigma$ est $\Sigma^{*}$, et l'inclusion
donne une équivalence entre revêtements étales, d'où une équivalence de
groupoïdes fondamentaux
\[
(14),(15)\qquad \Pi_1(\widetilde\Sigma) \xrightarrow{\ \approx\ } \Pi_1(\Sigma^{*}), \qquad
\Pi_1\Bigl(\coprod U_i\Bigr) \longrightarrow \Pi_1(\Sigma^{*}) .
\]
Les points de $U_i$ sont ce qu'on appelle aujourd'hui des \emph{points base
tangentiels}\footnote{La notion a été formalisée par Deligne (1989) pour la
droite projective moins trois points, comme base de l'action de Galois sur le
groupe fondamental ; ces pages n'en ont pas le nom.}. On dessine
$\widetilde\Sigma$ dans $\Sigma^{*}$ comme si l'on avait ôté trois disques
« infiniment petits » autour des trous, et l'on travaille avec des points de
$\widetilde\Sigma$ ou avec des parties simplement connexes de
$\widetilde\Sigma$. Le groupe $\Gamma_{\Sigma,J}$ agit sur toute la situation,
et « un aspect essentiel du travail à faire est de garder cette opération
constamment présente à l'esprit, et d'introduire des notations et des
constructions qui soient invariantes ».

\pagerange{49}{53}
\subsection*{Points, chemins et régions de base (pages 49 à 53)}

On introduit
\[
(16)\qquad Q_i = -R_i, \qquad R_i^{\omega} \in U_i\ (\omega \in \boldsymbol{\omega}), \qquad P_\varepsilon\ (\varepsilon \in \boldsymbol{\varepsilon}), \qquad S_i^{\varepsilon} \in U_i ,
\]
$Q_i$ le point opposé au trou, $P_\varepsilon$ le pôle de l'hémisphère
$\varepsilon$, $R_i^{\omega}$ la direction de l'équateur orientée par $\omega$
en $R_i$, $S_i^{\varepsilon}$ la direction du méridien de $R_i$ vers
l'hémisphère $\varepsilon$. Puis des chemins simples, par paquets
\[
(20)\qquad a_i^{\omega,\varepsilon}\ (12), \qquad b_i^{\omega}\ (6), \qquad c_i^{\varepsilon}\ (6), \qquad d_i^{\varepsilon}\ (6) :
\]
$a_i^{\omega\varepsilon} : S_i^{\varepsilon} \to R_i^{\omega}$ le quart de tour
sur $U_i$ ; $b_i^{\omega} : Q_i \to R^{\omega'}_{\omega' i}$ le long de
l'équateur, dans la direction $\omega$ ; $c_i^{\varepsilon} : P_\varepsilon
\to Q_i$ et $d_i^{\varepsilon} : P_\varepsilon \to S_i^{\varepsilon}$ le long
de méridiens\footnote{En (25) et (26), l'indice supérieur est écrit $\omega$ là
où (20) porte $\varepsilon$ ; on suit (20).}. Les paquets $b$, $c$, $d$ sont
des torseurs sous $\mathfrak{S}_J$, le paquet $a$ sous $\Gamma_{\Sigma,J}$ ;
$\tau$ échange les exposants $\varepsilon$ et fixe les $b$ (22).

\emph{Régions de base.} On découpe ensuite $\widetilde\Sigma$ par des
« incisions » qui relient les trois trous, de façon à obtenir des ouverts
simplement connexes aussi grands que possible :
\[
(27)\qquad \widetilde\Sigma_\varepsilon = \widetilde\Sigma \setminus \bigcup_{i} \operatorname{supp} d_i^{\varepsilon'}, \qquad
\widetilde\Sigma_i = \widetilde\Sigma \setminus \bigcup_{j \neq i} b_j ,
\]
$b_j$ étant l'arc d'équateur qui passe par $Q_j$ entre les deux trous
voisins\footnote{La page écrit $\operatorname{supp} d_i^{\varepsilon}$, et
$b_j = \operatorname{supp} c_j^{\omega} \cup \operatorname{supp} c_j^{\omega'}$.
Le dessin en marge montre les incisions issues du pôle opposé, et c'est la
seule lecture pour laquelle le « centre » $P_\varepsilon$ reste dans la région
; de même les incisions de $\widetilde\Sigma_i$ sont dessinées « le long de
l'équateur », donc faites des arcs $b$. Simple connexité : le complémentaire
est réunion des trois disques fictifs et d'un arbre qui les joint, ensemble
contractile.}. Ces régions, de « centres » $P_\varepsilon$ et $Q_i$, sont
connexes et simplement connexes, et l'on pose
\[
(28)\qquad \Pi_\varepsilon = \pi_1(\widetilde\Sigma; \widetilde\Sigma_\varepsilon) \simeq \pi_1(\Sigma^{*}, P_\varepsilon), \qquad
\Pi_i = \pi_1(\widetilde\Sigma; \widetilde\Sigma_i) \simeq \pi_1(\Sigma^{*}, Q_i) :
\]
le groupe fondamental basé en une région simplement connexe, canoniquement
isomorphe au groupe basé en n'importe lequel de ses points par des
« isomorphismes de transition » $\varphi$ (29). Les méridiens $c_i^{\varepsilon}$
donnent des isomorphismes $\Pi_\varepsilon \rightleftarrows \Pi_i$ (30), qui
relient ces cinq groupes de façon transitive, mais pas par un système
transitif : le long des cycles du graphe de l'octaèdre (deux pôles, trois
$Q_i$), on obtient des automorphismes intérieurs, et on les obtient
tous\footnote{Les trois demi-méridiens qui portent les $c$ découpent la sphère
en trois fuseaux contenant chacun un trou ; la sphère trouée se rétracte sur
ce graphe, dont les cycles engendrent donc le groupe fondamental. Le centre de
$\Pi_{0,3}$ étant trivial, on obtient bien tous les automorphismes intérieurs
(justification de l'édition).}. On note $\varphi_i$, $\varphi_\varepsilon$ les
repérages aux points $R_i^{\omega}$, $S_i^{\varepsilon}$ (31), le point base ne
faisant pas de doute dans les calculs.

\pagerange{53}{59}
\subsection*{Le groupe standard, et les épinglages (pages 53 à 59)}

On repère tout par le groupe standard
\[
(32)\qquad \Pi_{0,3} = \pi_1\bigl(U_{0,3}(\mathbb{C}), P^{+}\bigr) \subset
\widetilde\Pi_{0,3} = \pi_1\bigl(U_{0,3}(\mathbb{C}), \mathfrak{S}_3 \times \{1, \tau\}; P^{+}\bigr)
\]
— le groupe fondamental équivariant, formé des couples d'une symétrie et d'un
chemin de $P^{+}$ à son image, celui de la page 14 du dossier 141 —, avec
\[
\widetilde\Pi_{0,3} = \langle\, \rho, \sigma_\infty, \tilde\sigma_\infty \mid \rho^3 = \sigma_\infty^2 = \tilde\sigma_\infty^2 = 1,\ \sigma_\infty\rho\,\sigma_\infty^{-1} = \rho^{-1},\ (\sigma_\infty\tilde\sigma_\infty)^2 = 1 \,\rangle
\]
\[
\simeq GL(2, \mathbb{Z})/\{\pm 1\} = PGL_2(\mathbb{Z}) .
\]
La présentation est bien celle de $PGL_2(\mathbb{Z})$, produit amalgamé d'un
groupe diédral d'ordre $6$ et d'un groupe de Klein le long d'un
$\mathbb{Z}/2$, et l'identification est écrite ici de sa main\footnote{La
relation $\tilde\sigma_\infty(\sigma_\infty) = 1$ est lue avec doute ; on la
comprend comme la commutation de $\sigma_\infty$ et $\tilde\sigma_\infty$, ce
que confirme le crochet « i.e. $(\sigma_\infty\tilde\sigma_\infty)^2 = 1$ ».
Une version biffée donnait pour la partie orientée
$SL(2,\mathbb{Z})/\{\pm 1\}$. Le dossier 145 retrouve ce groupe comme groupe de
Coxeter $(2, 3, \infty)$ et y donne l'identification comme de l'édition ; elle
est en fait sur cette page-ci. La page suivante écrit
« $\widetilde\Pi_{0,3} \simeq SL(2,\mathbb{Z})$ » : il faut lire
$PGL_2(\mathbb{Z})$.}. Le « formulaire familier » (33) :
\[
\varepsilon_0 = \sigma_\infty\rho, \quad \varepsilon_1 = \rho(\varepsilon_0), \quad \varepsilon_\infty = \rho^2(\varepsilon_0), \qquad
l_i = \varepsilon_i^2, \qquad l_\infty l_1 l_0 = 1, \qquad \rho(l_i) = l_{\rho(i)},
\]
avec les valeurs de $\sigma_\infty$, $\tilde\sigma_\infty$, $\tau_\infty$ sur les
$l_i$ données dans les conventions\footnote{La page 54 donne
$\tau_\infty(l_\infty) = l_0^{-1}l_\infty^{-1}l_0$ (lecture douteuse) ; les
valeurs $\tau_\infty(l_0) = l_0^{-1}$, $\tau_\infty(l_1) = l_1^{-1}$, qui sont
sûres, imposent $\tau_\infty(l_\infty) = l_0 l_1 = l_0\, l_\infty^{-1}\,
l_0^{-1}$, comme l'écrit la page 23. Les $\varepsilon_i$ sont les éléments
paraboliques de $\widetilde\Pi_{0,3}$ que le dossier 145 appelle racines
quasi-carrées ; la lecture $\varepsilon_0 = \sigma_\infty\rho$ est douteuse,
mais elle s'accorde avec celle de 145, $\varepsilon_1 = \sigma_0\rho$, par
conjugaison par $\rho$.}.

\emph{Épinglages.} Un \emph{repère} de la situation est un élément
$r = (i, \omega, \varpi) \in J \times \boldsymbol{\omega} \times
\boldsymbol{\varpi}$, qui définit $\varepsilon = \omega \wedge \varpi$. À
chaque repère correspond l'unique isomorphisme (conforme ou anticonforme)
$u_r : (U_{0,3}(\mathbb{C}), P^{+}) \to (\Sigma^{*}, P_\varepsilon)$ qui envoie
le repère standard sur $r$, et l'\emph{isomorphisme d'épinglage}
\[
(34),(36)\qquad \psi_r = \psi_\varepsilon^{i,\omega} = \pi_1(u_r; P^{+}) : \Pi_{0,3} \xrightarrow{\ \sim\ } \Pi_\varepsilon ,
\]
et de même $\psi_i^{\omega,\varepsilon} : \Pi_{0,3} \xrightarrow{\sim} \Pi_i$,
compatibles par les isomorphismes (30). Pour $P_\varepsilon$ fixé il y a six
épinglages, pour $Q_i$ fixé quatre ; trois et deux si $\Sigma$ est
orientée\footnote{La page écrit « trois resp. », le second nombre étant
remplacé d'un premier « douze » biffé ; avec $\varpi$ fixé, $\varepsilon$ est
déterminé par $\omega$, d'où $3$ choix en $P_\varepsilon$ et $2$ en $Q_i$.}.
Le stabilisateur de $P_\varepsilon$ dans $\Gamma_{\Sigma,J}$,
\[
\Gamma_{P_\varepsilon} = \{(g, \tau^{\alpha}) \mid \alpha = 0 \text{ si } g \text{ est pair},\ \alpha = 1 \text{ sinon}\},
\]
le relèvement canonique de $\mathfrak{S}_3$ que le dossier 145 note
$\Gamma_P$ (sa page 20), agit simplement transitivement sur les repères de
même $\varepsilon$, et
\[
(39)\qquad \psi_{r g_0} = \psi_r \circ \pi_1(g_0; P^{+}), \qquad
\psi_\varepsilon^{\omega i, \omega} = \psi_\varepsilon^{i,\omega} \circ \rho, \qquad
\psi_\varepsilon^{\omega' i, \omega'} = \psi_\varepsilon^{i,\omega} \circ \tilde\sigma_\infty ,
\]
$\rho$ s'interprétant comme la rotation d'un tiers de tour autour de l'axe
polaire, $\tilde\sigma_i$ comme la symétrie par rapport au méridien de $Q_i$.
En $Q_i$, le stabilisateur de $Q_0$ dans le groupe standard est un groupe de
Klein engendré par $\sigma_\infty$ et $\tau$, et (pages 58–59)
\[
\psi_i^{\omega,\varepsilon'} = \psi_i^{\omega,\varepsilon} \circ \tau_\infty, \qquad
\psi_i^{\omega',\varepsilon'} = \psi_i^{\omega,\varepsilon} \circ \sigma_\infty, \qquad
\psi_i^{\omega',\varepsilon} = \psi_i^{\omega,\varepsilon} \circ \tilde\sigma_\infty ,
\]
la troisième relation étant composée des deux premières
($\tilde\sigma_\infty = \sigma_\infty\tau_\infty$)\footnote{Pour cette
comparaison, $\pi_1(U_{0,3}; Q_\infty)$ est identifié à $\Pi_{0,3}$ par le
chemin $d_\infty^{+}$ de $P^{+}$ à $Q_\infty$, alors qu'en (25)–(26) c'est $c$
qui mène aux $Q_i$ ; les numéros (39), (40) sont repris deux fois.}.

\pagerange{59}{63}
\subsection*{Les régions-drapeaux, et les points proches des trous (pages 59 à 63)}

« Dans tous ces calculs, suivant sans doute d'anciennes habitudes, j'ai
privilégié les $P_\varepsilon$ et les $Q_i$ comme points base — alors que pour
le travail que j'ai en vue, ce sont les points $R_i^{\omega}$ qui sont les plus
cruciaux », parce que les calculs en ces points expriment directement la
monodromie autour des trous. Il introduit donc, pour chaque repère $r$, une
région associée au triangle-drapeau $D_r = P_\varepsilon Q_i R_{\omega i}$ de
la subdivision barycentrique de la sphère par l'équateur et les trois
trous — douze triangles —, tronqué sur $\widetilde\Sigma$ en un quadrilatère
$\widetilde D_r$, et sa réunion avec trois de ses symétriques\footnote{La page écrit le dernier terme $\tilde\sigma_{\omega i}(\Sigma_r)$ ;
son contour n'est lu qu'en partie, et l'on n'en donne pas la liste des
sommets. Celui de la région standard, page 61, est
$P_+ Q_1 R_0^- S_0^+ R_0^+ S_0^- P_- Q_\infty R_1^- S_1^+ P_+$.} :
\[
(41)\qquad \widetilde\Sigma_r = \widetilde D_r \cup \tau(\widetilde D_r) \cup \tilde\sigma_i(\widetilde D_r) \cup \tilde\sigma_{\omega i}(\widetilde D_r) .
\]
Ces douze
régions sont simplement connexes et permutées simplement transitivement par
$\Gamma_{\Sigma,J}$ ; d'où des groupes $\Pi_r = \pi_1(\widetilde\Sigma;
\widetilde\Sigma_r)$ reliés par un système transitif d'isomorphismes, et des
épinglages $\Psi_r : \Pi_{0,3} \to \Pi_r$ (42). Un point $R_i^{\omega}$ est
intérieur à exactement deux régions $\widetilde\Sigma_r$ ; si $\Sigma$ est
orientée par $\varpi_+$, l'une est privilégiée, et l'on obtient
\[
(43)\qquad \Psi_i^{\omega} : \Pi_{0,3} \xrightarrow{\ \sim\ } \pi_1(\widetilde\Sigma, R_i^{\omega}) ,
\]
compatible aux repérages en $Q_i$ et $P_\omega$ par les chemins $b$ et $c$
(44). Le passage d'un côté à l'autre du trou, par le chemin $\tilde b_i^{\omega}$
qui passe par $Q_i$, s'écrit
\[
(45)\qquad \tilde b_i^{\omega}\bigl(\Psi_i^{\omega}(x)\bigr) = \Psi_i^{\omega'}\bigl(\sigma_\infty(x)\bigr) ,
\]
la forme même des formules de transport de la suite II (comparer (6) et (42)
de la page 30). La suite s'arrête là, au haut de la page 63.

\pagerange{65}{103}
\section*{IV. Généralités sur les actions de groupes de Galois sur les complétés profinis de groupoïdes fondamentaux (pages 65 à 103)}

La page 65 est une couverture de sa main, avec au crayon « (38 pages) » ; la
suite est paginée de 1 à 18, puis continue.

\pagerange{66}{69}
\subsection*{Le cadre (pages 66 à 69)}

Soit $K$ un sous-corps de $\mathbb{C}$ — surtout de type fini sur
$\mathbb{Q}$, la clôture algébrique $\overline{\mathbb{Q}}$ en étant un cas
limite, « pour ne pas être gêné par les racines de 1 » —, et $X$ un schéma de
type fini sur $K$, qu'on regarde comme l'espace $X(\mathbb{C})$ muni d'une
structure supplémentaire ; on traite les homéomorphismes universels comme des
isomorphismes. Le groupoïde fondamental topologique $\Pi_1(X)$ se plonge dans
son complété profini $\widehat\Pi_1(X)$, bijectivement sur les objets, et
\[
(3)\qquad \pi_1(X_{\mathrm{sch}}, x) \simeq \widehat\pi_1(X, x)
\]
(théorème d'existence de Riemann, SGA~1, exposé XII) ; les ensembles de
flèches $\mathrm{Hom}_{\widehat\Pi_1}(x, y)$ sont les bitorseurs déduits de
ceux de $\Pi_1$ par extension des groupes. Comme $\widehat\Pi_1$ a une
définition purement schématique, le groupe
\[
(5)\qquad \Gamma_K = \mathrm{Aut}(\mathbb{C}/K)
\]
agit sur lui, par ${}^{u}\lambda : {}^{u}x \to {}^{u}y$ pour $\lambda : x \to y$,
de façon compatible à la composition (6) ; une flèche de $\Pi_1$ n'a en
général qu'une image dans $\widehat\Pi_1$. Sur le sous-groupoïde induit sur
les points algébriques $X(\overline{K})$, l'action se fait à travers
$\mathrm{Gal}(\overline{K}/K)$ (7)\footnote{La page écrit ici $\Pi_1$ sans
chapeau ; il s'agit de $\widehat\Pi_1(X_{\overline{K}})$, comme le dit la
phrase précédente. Que l'action de $\mathrm{Aut}(\mathbb{C}/\overline{K})$
soit triviale sur ce sous-groupoïde tient à l'invariance du groupoïde
fondamental étale par extension de corps algébriquement clos en
caractéristique nulle.}.

Pour des points rationnels $x, y \in X(K)$, fixés par $\Gamma_K$, et un chemin
$l : x \to y$ de $\Pi_1$, on écrit
\[
(8),(9)\qquad {}^{u}l = g_l(u)\, l = l\, f_l(u),
\]
\[
f_l(u) = l^{-1}\,{}^{u}l \in \widehat\pi_1(X, x), \quad g_l(u) = {}^{u}l\, l^{-1} = l\bigl(f_l(u)\bigr) \in \widehat\pi_1(X, y),
\]
et l'on décrit l'action de $u$ par les $f_l(u)$, ou les $g_l(u)$, pour $l$
parcourant des générateurs — en nombre fini sur une partie finie $I$ de
$X(K)$ : un $x_0 \in I$, un chemin de $x_0$ à chaque autre point, des
générateurs de $\pi_1(X, x_0)$. Il présume que pour une courbe anabélienne la
connaissance de l'action de $u$ sur $\widehat\Pi_1(X, I)$ détermine $u$.

\pagerange{70}{77}
\subsection*{Multiplicativité, régions de base, cocycles (pages 70 à 77)}

Pour $x \xrightarrow{l'} y \xrightarrow{l} z$,
\[
(10)\qquad g_{ll'}(u) = g_l(u)\, l\bigl(g_{l'}(u)\bigr), \qquad
(11)\qquad f_{ll'}(u) = l'^{-1}\bigl(f_l(u)\bigr)\, f_{l'}(u) ,
\]
et (12) itère (10)\footnote{Vérifié : ${}^{u}(ll') = g_l\, l\, g_{l'}\, l' =
g_l\,(l g_{l'} l^{-1})\,ll'$, et de même pour $f$. Un premier (12), qui
parlait de 1-cocycles en $u$, est barré ; il reviendra en (29).}. Ces formules
« d'allure peu sympathique » se simplifient si l'on remplace les chemins par
des \emph{régions de base} : soit $p : T \to X$ continue, $T$ simplement
connexe, et $\pi = \pi_1(X; T)$, qui s'identifie à chaque $\pi_1(X, p(x))$ par
$\varphi_x$ (15) ; deux points $x, y \in T$ donnent une flèche canonique
$l_{y,x} : p(x) \to p(y)$ (16). Pour $x, y \in T_K = p^{-1}(X(K))$, $f$ et $g$
deviennent un même élément de $\widehat\pi$,
\[
(17)\qquad g^{T}_{y,x}(u) = \varphi_y^{-1}\bigl(g_{l_{y,x}}(u)\bigr) = \varphi_x^{-1}\bigl(f_{l_{y,x}}(u)\bigr),
\]
et (10), (11) se réduisent à la relation de Chasles
\[
(18)\qquad g^{T}_{z,y}(u)\, g^{T}_{y,x}(u) = g^{T}_{z,x}(u) \qquad (x, y, z \in T_K) ;
\]
« quant aux formules de passage (9), elles disparaissent ! »\footnote{Vérifié :
$\varphi_y = l_{y,x} \circ \varphi_x$ comme transports, d'où l'égalité de
(17) ; (18) découle de (10). La marge appelle $T$ « la conjugaison » dans le
cas le plus intéressant : vraisemblablement $T$ un demi-plan, image d'un
domaine fondamental de la conjugaison complexe ; c'est une lecture.}. Les
invariants $f_l$, $g_l$ sont le cas $T = [0, 1]$ (20), et, pour $T \to T'$ au-dessus
de $X$, $g^{T'} = \pi_1(f)(g^{T})$ (23).

Soit alors $X = \bigcup_\alpha T_\alpha$, les $T_\alpha$ simplement connexes,
d'intérieurs recouvrant $X$, ordonnés par inclusion de sorte que $T_\alpha \cap
T_\beta = \bigcup_{\gamma \leq \alpha, \beta} T_\gamma$ (28), avec les groupes
$\pi_\alpha = \pi_1(X, T_\alpha)$ et, pour $T_\alpha \subset T_\beta$,
l'isomorphisme $\varphi_{\beta\alpha} : \pi_\alpha \to \pi_\beta$, de sorte que
$g^{\beta}_{y,x} = \varphi_{\beta\alpha}(g^{\alpha}_{y,x})$ (27)\footnote{La page
écrit $\varphi_{\beta\alpha} : \pi_\beta \to \pi_\alpha$ en (26), puis
l'applique à $g^{\alpha}$ en (27) ; c'est un isomorphisme, on l'oriente dans
le sens où il s'applique.}. La connaissance des $g^{\alpha}_{y,x}(u)$
détermine-t-elle tous les $g_l(u)$ ? L'argument de la page subdivise un chemin
en petits chemins entre points « voisins » ; il suppose des points
$K$-rationnels partout, ce qui n'est pas le cas en général. L'énoncé juste est
celui que la page donne ensuite, « un peu plus général, et peut-être plus
intéressant » : soit $I \subset X(K)$ fini, $I_\alpha = I \cap T_\alpha$ ; si
$I$ rencontre chaque composante connexe par arcs de chaque intersection finie
des $T_\alpha$, le théorème de van Kampen pour les groupoïdes montre que
$\Pi_1(X, I)$ est engendré par les flèches $l^{\alpha}_{y,x}$, $x, y \in
I_\alpha$, avec pour relations les identifications le long des inclusions ;
la connaissance des $g^{\alpha}_{y,x}(u)$ pour $x, y \in I_\alpha$ détermine
alors \emph{tous} les $g_l(u)$, $l$ chemin de $\Pi_1(X, I)$\footnote{Le
théorème de van Kampen pour les groupoïdes est de R. Brown (1967) ; la page ne
le cite pas. Qu'il faille une hypothèse sur $I$ : pour une courbe de genre
$\geq 2$ sur un corps de nombres, $X(K)$ est fini (Faltings, 1983), et
l'argument de subdivision ne peut s'appliquer. La page parle aussi d'un
« argument de descente » pour des $T_\alpha$ fermés, sans le donner.}.

\emph{Variation de $u$.} De ${}^{uv}l = {}^{u}({}^{v}l)$ (28) on tire
\[
(29),(30)\qquad g_l(uv) = {}^{u}g_l(v)\, g_l(u), \qquad f_l(uv) = f_l(u)\, {}^{u}f_l(v) :
\]
pour $l$ fixé, $f_l$ et $g_l^{-1}$ sont des 1-cocycles (homomorphismes
croisés) de $\Gamma_K$ à valeurs dans $\widehat\pi_1(X, x)$, resp.
$\widehat\pi_1(X, y)$, déduits l'un de l'autre par le transport le long de $l$
; ce sont les cocycles du torseur $\mathrm{Hom}_{\widehat\Pi_1}(x, y)$, dont $l$
est un « élément origine »\footnote{Vérifié. La page dit $f$ cocycle « droit »
et $g$ cocycle « à gauche » ; dans les conventions usuelles, $f_l$ est un
1-cocycle ordinaire pour l'action de $\Gamma_K$ sur $\widehat\pi_1(X, x)$.}.

\pagerange{78}{82}
\subsection*{Classes de cohomologie, et Mordell–Weil (pages 78 à 82)}

Si l'on remplace $l$ par $l' = la = bl$,
\[
(31)\qquad f_{l'}(u) = a^{-1}\, f_l(u)\, {}^{u}a, \qquad g_{l'}(u) = {}^{u}b\, g_l(u)\, b^{-1} :
\]
les cocycles changent en des cocycles cohomologues\footnote{Vérifié.}. Deux
points rationnels $x, y$ définissent donc des classes
\[
(32)\qquad f_{y,x} \in H^1\bigl(\Gamma_K, \widehat\pi_1(X, x)\bigr), \qquad g_{y,x} \in H^1\bigl(\Gamma_K, \widehat\pi_1(X, y)\bigr)
\]
dans les ensembles de cohomologie non abélienne, qui se déterminent
mutuellement, de même que $f_{x,y}$ et $f_{y,x}$ ($g_{l^{-1}} = f_l^{-1}$,
(34))\footnote{La page écrit $\widehat\pi_1(Y, y)$ en deux endroits pour
$\widehat\pi_1(X, y)$.}. Toute cette structure se décrit par le bitorseur
$\mathrm{Hom}_{\widehat\Pi_1}(x, y)$ et l'action de Galois sur « tout ce
fourbi ».

Quand $X$ se plonge dans une extension d'une variété abélienne par un tore
— par exemple une courbe anabélienne dans sa jacobienne généralisée — et que
$K$ est de type fini sur $\mathbb{Q}$, le théorème de Mordell–Weil entraîne
qu'\emph{un point $x \in X(K)$ est entièrement déterminé par la classe
$f_{x_0,x}$}, $x_0$ fixé, et même par son image abélianisée
\[
(34)\qquad (f_{x_0, x})_{\mathrm{ab}} \in H^1\bigl(\Gamma_K, \widehat\pi_1(X, x_0)_{\mathrm{ab}}\bigr) .
\]
L'énoncé est juste\footnote{En
effet, par l'application d'Albanese, la classe abélianisée est l'image du point
par l'application de Kummer $A(K) \to H^1(\Gamma_K, \widehat{T}(A))$ de la
variété semi-abélienne, injective parce que $A(K)$ est de type fini
(Mordell–Weil ; Lang–Néron pour $K$ de type fini). C'est l'injectivité de ce
qu'on appelle aujourd'hui l'application de Kummer, ou des sections, de la
géométrie anabélienne ; la surjectivité est l'objet de la conjecture des
sections, que Grothendieck formulera dans sa lettre à Faltings de 1983.}.

Déterminer « explicitement » un tel cocycle, en termes des propriétés
arithmétiques du point, « c'est là un travail fascinant en perspective, que je
n'ai pratiquement pas encore commencé ! ». Mais que veut dire
« explicitement », le groupe $\Gamma_K$ étant lui-même loin d'être connu ? Il
faut repérer $u$ par des invariants $f_{l_j}(u)$ attachés à des générateurs
d'un groupoïde $\Pi_1(X, I)$. Or pour une courbe elliptique, ou pour
$\mathbb{G}_m$, l'action de $\Gamma_K$ sur $\widehat\Pi_1(X_{\overline{K}})$
n'est pas fidèle : il faut des paramètres de repérage pris sur d'autres
variétés, la plus naturelle étant
\[
(35)\qquad U_{0,3} = \mathbb{P}^1 \setminus \{0, 1, \infty\} .
\]

\pagerange{82}{89}
\subsection*{Fonctorialité, symétries, réduction à $U_{0,3}$ (pages 82 à 89)}

Un morphisme $f : X \to Y$ de $K$-schémas donne $\widehat f :
\widehat\Pi_1(X) \to \widehat\Pi_1(Y)$, qui commute à l'action de Galois (39),
d'où $\widehat f(f_l(u)) = f_{\widehat f(l)}(u)$ et de même pour $g$ (40). Le
cas extrême est celui d'un groupe $\mathfrak{G}$ de $K$-automorphismes de $X$ :
$\widehat\Pi_1(X)$ devient un groupoïde à opérateurs $\mathfrak{G} \times
\Gamma_K$. Avec un point base $x_0$, l'action de $\Gamma_K$ sur
$\widehat\Pi_1(X, I)$ est connue par son action sur $\widehat\pi_1(X, x_0)$ et
sur les torseurs $\mathrm{Hom}(x_0, x_i)$, soit par les données (42), (43) ;
mais le choix d'un point base introduit une dissymétrie « qui semble
artificielle, et à la longue je crois impraticable », surtout en présence de
$\mathfrak{G}$, « alors que pratiquement jamais on ne peut trouver un
$x \in X(K)$ fixe sous $\mathfrak{G}$ ». C'est exactement la difficulté que
les suites II et III contournaient par des systèmes symétriques de points et
de régions de base.

\emph{Réduction à $U_{0,3}$.} Soit $X$ quasi-projective lisse
géométriquement connexe. Par un théorème de type Lefschetz, $X$ contient une
courbe $C$ dont l'inclusion est surjective sur les groupes fondamentaux ; par
le théorème de Belyi, un ouvert de $C$ est un revêtement fini étale de
$U_{0,3}$ :
\[
(44)\qquad
\begin{array}{ccc}
C & \overset{i}{\hookrightarrow} & X\\
{\scriptstyle f}\downarrow & & \\
U_{0,3} & &
\end{array}
\qquad i \text{ surjectif sur les } \pi_1,\quad f \text{ fini étale},\quad I \subset C .
\]
Alors $\widehat\Pi_1(X, I)$ est un quotient de $\widehat\Pi_1(C, I)$, et
$\widehat\Pi_1(C, I) \to \widehat\Pi_1(U_{0,3}, f(I))$ est fidèle,
« injective » sur les ensembles de flèches. Les deux flèches commutant à
Galois, \emph{l'action de $\Gamma_K$ sur $\widehat\Pi_1(U_{0,3}, I_0)$, $I_0 =
f(I)$, détermine son action sur $\widehat\Pi_1(X, I)$}\footnote{La page passe
par le saturé $I' = f^{-1}f(I)$, puis remarque (page 88) qu'il n'est pas
rationnel en général et qu'on n'en a pas besoin ; c'est la version sans
saturé qu'on donne ici. Belyi demande que $C$ soit définie sur un corps de
nombres : la réduction vaut pour $K$ algébrique sur $\mathbb{Q}$, quitte à
remplacer $\Gamma_K$ par un sous-groupe ouvert (pages 88–89) sur lequel $C$,
$f$ et $i$ sont définis. Dans les deux formules de la page 87 où l'on attend
$C$, la page écrit $U_{0,3}$. Le théorème de Belyi est de 1979 ; sa mention
date au moins cette suite.}.

Inversement, on cherchera des conditions nécessaires et suffisantes pour
qu'un automorphisme $\Theta$ de $\widehat\Pi_1(U_{0,3}; I_0)$ fixant les
sommets provienne d'un $u \in \Gamma_K$. Première condition : pour tout
revêtement $f : C \to U_{0,3}$ défini sur $K$, avec $I = f^{-1}(I_0) \subset
C(K)$, $\Theta$ doit se relever en une action sur $\widehat\Pi_1(C, I)$ qui
soit l'identité sur les sommets ; quand $I \to I_0$ n'est pas surjectif, ou
les points de $I$ pas rationnels, on demande la même chose pour une puissance
$\Theta^n$ convenable.

\pagerange{90}{97}
\subsection*{Un lemme de relèvement (pages 90 à 97)}

À ces conditions de « prolongeabilité » s'en ajoute une : pour un morphisme
$C \to X$, que l'action passe au quotient, c'est-à-dire que l'action sur
$\widehat\pi_1(C, x_0)$ stabilise le noyau de $\widehat\pi_1(C, x_0) \to
\widehat\pi_1(X, x_0)$. Pour exprimer le relèvement, la page énonce :

\textbf{Lemme.} \emph{Soit $\widetilde{\mathcal{C}} \to \mathcal{C}$ un
revêtement de groupoïdes, $\mathcal{C}$ connexe, $x_0$ un sommet de
$\mathcal{C}$, $x'_0$ un sommet au-dessus, $\Pi_0 = \pi_1(\mathcal{C}, x_0)$
et $\Pi'_0 \subset \Pi_0$ l'image de $\pi_1(\widetilde{\mathcal{C}}, x'_0)$.
Soit $\Theta$ un automorphisme de $\mathcal{C}$ fixant $x_0$. Alors $\Theta$ se
relève en un automorphisme de $\widetilde{\mathcal{C}}$ fixant $x'_0$ si et
seulement si $\Theta(\Pi'_0) = \Pi'_0$ ; le relèvement est alors unique, et
sur les fibres il est donné par transport de structure à partir de}
\[
(45)\qquad \mathrm{Ob}\,\widetilde{\mathcal{C}}_{x_i} = H_{x_i, x_0} \times^{\Pi_0} \Pi_0/\Pi'_0, \qquad H_{x_i, x_0} = \mathrm{Hom}_{\mathcal{C}}(x_0, x_i) .
\]
\emph{Sans exiger que $x'_0$ soit fixe, $\Theta$ se relève si et seulement si
$\Theta(\Pi'_0)$ est conjugué de $\Pi'_0$ dans $\Pi_0$, et les relèvements
correspondent aux éléments de}
\[
(49)\qquad \mathrm{Transp}_{\Pi_0}\bigl(\Pi'_0, \Theta(\Pi'_0)\bigr)/\Pi'_0, \qquad \mathrm{Transp} = \{g \in \Pi_0 \mid g\,\Pi'_0\,g^{-1} = \Theta(\Pi'_0)\} .
\]
\emph{Pour un sous-groupoïde plein $\mathcal{C}' \subset
\widetilde{\mathcal{C}}$, équivalent à un revêtement de $\mathcal{C}$, il faut
de plus que $\mathrm{Ob}\,\mathcal{C}'$ soit stable par le relèvement.}

En effet un relèvement sur la fibre $E_0 = \Pi_0/\Pi'_0$ est une bijection
$\widetilde\Theta_0$ vérifiant $\widetilde\Theta_0(g x') =
\Theta(g)\widetilde\Theta_0(x')$ (46) ; elle est déterminée par l'image $g\Pi'_0$
du point base, et la compatibilité des stabilisateurs impose
$\Theta(\Pi'_0) = g\Pi'_0 g^{-1}$ ; $g$ et $g\alpha$, $\alpha \in \Pi'_0$,
donnent le même relèvement\footnote{Démonstration de l'édition, sur les
indications des pages 91 à 93, qui sont très corrigées. L'ensemble (49), s'il
est non vide, est en bijection avec $N_{\Pi_0}(\Pi'_0)/\Pi'_0$, le groupe des
automorphismes du revêtement.}.

Résumé (page 94) : pour qu'un automorphisme $\Theta$ de $\widehat\Pi_1(X, I)$,
$I \subset X(K)$, provienne d'un $u \in \Gamma_K$, il faut que
\begin{itemize}
\item[a)] pour tout revêtement fini étale $X'$ de $X$ défini sur $K$, $\Theta$
se relève en une action sur $\widehat\Pi_1(X', I')$, $I' = X'|_I$, qui soit
l'identité sur $I' \cap X'(K)$ ;
\item[b)] pour tout morphisme $f : X' \to Y$ défini sur $K$ et tout $x' \in
I' \cap X'(K)$, le noyau de $\widehat\pi_1(X', x') \to \widehat\pi_1(Y, f(x'))$
soit stable par le relèvement — et, sans supposer $x'$ rationnel, par
$\Theta^{\nu}$ pour $\nu$ tel que $\Theta^{\nu}x' = x'$.
\end{itemize}
« Je ne pense pas trop exagérer en conjecturant » que ces conditions
nécessaires sont aussi suffisantes, quand $X$ est une courbe anabélienne et
$K$ de type fini sur $\mathbb{Q}$. Quant à l'unicité de $u$, il l'attend quand
$K$ est algébrique sur $\mathbb{Q}$, et, en degré de transcendance positif,
quand $K$ est fini sur le corps des modules de $X$\footnote{La page 95 est en
partie illisible. L'unicité, c'est la fidélité de l'action extérieure de
Galois ; elle est connue pour $U_{0,3}$ sur $\overline{\mathbb{Q}}$ (Belyi),
et pour les courbes hyperboliques sur les corps de nombres (Matsumoto, 1996,
pour les affines ; Hoshi et Mochizuki, 2011, pour les propres). La suffisance
— une caractérisation de l'image de Galois par des conditions de cette sorte
— est, à notre connaissance, ouverte.}.

Pages 96–97 : le cas d'un revêtement $\widetilde{\mathcal{C}}$ non connexe,
venant d'un $X' \to X$ seulement $K$-connexe : sur la fibre $E_0$, le groupe
engendré par $\Gamma_K$ et $\Pi_0$ agit transitivement, sans que $\Pi_0$ le
fasse. L'application qui à $x' \in E_0$ associe son stabilisateur induit
\[
(51)\qquad \pi_0(E_0) = E_0/\Pi_0 \longrightarrow \mathrm{Ssgr}(\Pi_0)/\Pi_0
\]
vers les classes de conjugaison de sous-groupes, et une condition nécessaire
évidente est que l'image soit stable par $\Theta$ ; un relèvement est alors
déterminé par une bijection de $\pi_0(E_0)$ compatible à (51) et par l'image
d'un point dans chaque composante\footnote{Page très corrigée ; l'ordre de
lecture des ajouts est incertain. La page 109 reprend ce diagramme, avec
$\Sigma = \Gamma \cdot \pi$, $\Sigma/\pi = \Gamma$.}.

\pagerange{98}{103}
\subsection*{La condition des lacets, et les conditions de stabilité (pages 98 à 103)}

Il reste une condition, « plus importante » : la compatibilité avec la
« structure à l'infini » — pour une courbe, la donnée des classes de
conjugaison des sous-groupes d'inertie aux trous — qui s'exprime
« particulièrement élégamment » dans le langage des groupoïdes. Les conditions
a), b), au contraire, ne s'expriment pas directement sur les paramètres
($\alpha, \beta, \zeta, \eta, \ldots$) qui repèrent un $u \in
\mathrm{Gal}(\overline{\mathbb{Q}}/\mathbb{Q})$ par son action sur
$\Pi_{0,3}$. Il annonce, à propos des $M_{g,\nu}$ et des groupes de
Teichmüller, une conjecture « intéressante et plausible » en termes de
propriétés des groupes profinis tels que $\widehat\Pi_{0,3}$, liées à
$M_{0,4}$, $M_{0,5}$\footnote{La page 99 est presque entièrement illisible ;
on ne garde que ces mots. On pense à la relation pentagonale, sur $M_{0,5}$,
que l'\emph{Esquisse d'un programme} placera au cœur du programme (« les deux
premiers étages de la tour de Teichmüller ») et que Drinfeld prendra pour
troisième relation de $\widehat{GT}$ ; le rapprochement est de l'édition.}.
Puis il récapitule, pour un automorphisme $\Theta$ de $\widehat\Pi_1(X, I)$,
$I \subset X(\overline{K})$, $I \cap X(K) \neq \emptyset$ — surtout pour
$X = U_{0,3}$ :
\begin{itemize}
\item[1)] la \emph{condition des lacets}, « pour mémoire — c'est sûrement la
plus importante de toutes, et de très loin ! » : $\Theta$ envoie chaque lacet
autour d'un trou sur un conjugué de sa puissance $p$-ième, avec un même
multiplicateur $p$ ;
\item[2)] la stabilité pour les revêtements étales $X'$ de $X$ définis sur
$K$ : $\Theta$ se relève à $\widehat\Pi_1(X', I')$, de façon unique si $X'$
est $K$-connexe et $I' \cap X'(K) \neq \emptyset$ — ce qui s'explicite en une
condition de stabilité de certains sous-groupes ouverts de $\pi_1(X, x)$ ;
\item[3)] pour les morphismes $f : X' \to Y$ et $x' \in X'(K) \cap I'$, la
stabilité du noyau de $\widehat\pi_1(X', x') \to \widehat\pi_1(Y, f(x'))$ ;
\item[4)] pour deux revêtements $X'_1$, $X'_2$ de $X$ et un morphisme $f :
X'_1 \to X'_2$, pas nécessairement au-dessus de $X$, la compatibilité de
$\widehat\Pi_1(X'_1, J) \to \widehat\Pi_1(X'_2, I'_2)$, $J = I'_1 \cap
f^{-1}(I'_2)$, aux relèvements de $\Theta$.
\end{itemize}
Une condition 5° réunissant 3° et 4° serait facile à écrire, dit la page 101,
mais il ne sait pas s'il en aura besoin\footnote{Le haut de la page 101 est
très incertain ; le petit diagramme $X'_1 \to X'_2$, $f_1, f_2 : X'_i \to Y$ y
est trop raturé pour être rendu.}.

\emph{Adjacence (pages 102–103).} Soit $T \subset X(\mathbb{C})$ localement
fermée, simplement connexe, et $X' \to X$ fini étale au voisinage de $T$ ;
$X'|_T$ est réunion de copies de $T$, et deux points au-dessus de $T$ sont
\emph{$T$-adjacents} s'ils sont sur la même copie. Galois respecte
l'adjacence, et l'on définit $g^T_{x',y'}(u)$ par
${}^{u}l^T_{y',x'} = g^T(u)\, l^T_{u(y'),u(x')}$ (53), avec
$\widehat f(g^T_{y',x'}(u)) = g^T_{y,x}(u)$ (54). D'où une condition nécessaire :
\[
g^{T}_{y,x}(u) \in \bigcap_{\alpha \in \pi_0(X'|_T)} \mathrm{Im}\,\widehat\varphi_\alpha \qquad \text{pour tout } u ,
\]
$\widehat\varphi_\alpha : \widehat\pi_1(X', T_\alpha) \hookrightarrow
\widehat\pi_1(X, T)$ les inclusions des composantes $T_\alpha$, dont les images
sont des conjugués d'un même sous-groupe, sous la seule hypothèse — plus faible
que la rationalité — que $u(x'_\alpha)$ et $u(y'_\alpha)$ restent
$T$-adjacents\footnote{Ces deux pages sont d'une lecture très partielle ;
on ne retient que les formules et la conclusion. La marge identifie
$\pi_0(X'|_T)$ à la fibre de $X'$ au-dessus de $T$, sur laquelle agit
$\pi_1(X, T)$.}. La suite s'arrête là.

\pagerange{104}{121}
\section*{V. Groupoïde fondamental d'une carte trouée (pages 104 à 121)}

La page 104 est une feuille de garde de sa main, avec ce titre — sous-titré
« préliminaires à l'action de Galois » — et « (14 pages) ». Le titre reprend
la « carte trouée » dont la page 10 du dossier 141 commençait la définition :
une carte dont les sommets, les milieux d'arêtes et les centres des faces sont
percés.

\pagerange{105}{112}
\subsection*{Esquisses (pages 105 à 112)}

\emph{Présentations de $\Pi_1(U_{0,3})$ (page 105).} Quatre figures présentent
le groupoïde fondamental de la sphère privée de $0, 1, \infty$ sur divers
ensembles de points base, choisis sur de petits cercles autour des trous :
« 6 sommets, 9 flèches, 2 relations » ; « 4 sommets, 6 flèches, 1 relation »
(trois cas) ; « 2 sommets, trois flèches, 0 relation » (trois cas) ;
« 1 sommet, 2 flèches, 0 relation » (trois cas). Dans chaque cas le nombre
$\text{flèches} - \text{sommets} + 1 - \text{relations}$ vaut $2$, le rang du
groupe libre $\Pi_{0,3}$\footnote{Vérification de l'édition. Une cinquième
figure, « 1 cas », est biffée.}.

\emph{La sphère à quatre trous (page 107).} Les points $0, \infty$ aux pôles,
$1, -1$ (et $\pm i$) sur l'équateur ; les symétries $z \mapsto -z$,
$z \mapsto 1/z$, qui engendrent un groupe de Klein $\mathbb{F}_2^{\,2}$ ; la
suite exacte $1 \to \mathbb{F}_2^{\,2} \to \mathfrak{A}_4 \to \mathbb{Z}/3 \to 1$ ;
la relation $l_\infty\, l_1\, l_0\, l_{-1} = 1$ entre les quatre lacets\footnote{La page ne dit pas ce
que $\mathfrak{A}_4$ y fait ; la suite exacte est celle du groupe alterné,
exacte, et les deux involutions engendrent bien un groupe de Klein. On ne
reconstruit pas le reste.}.

\emph{Le modèle algébrique (page 110).} La sphère de III se réalise comme
$\Sigma_J = \mathbb{P}(V(J))$, $V(J) = \ker(\mathcal{O}^J \to \mathcal{O})$ ; le
fibré tangent est $\mathcal{O}(2)$ tordu par le $\mathbb{Z}/2$-torseur
$\boldsymbol{\omega}$ des orientations (« attention signe ! »), les trous
$R_i$, les points $Q_i$, $Q_i^{\omega}$ et les pôles $P_\omega$ — ceux-ci
définis à l'aide d'une racine cubique primitive de l'unité $\zeta$, comme
droites propres de la permutation circulaire — étant donnés par des vecteurs
explicites de $V(J)$\footnote{Formules en partie surchargées ; on n'en
reproduit pas le détail. Que les pôles demandent $\zeta$ est le fait
arithmétique notable : ils ne sont pas rationnels sur $\mathbb{Q}$, mais sur
$\mathbb{Q}(\zeta)$, et la conjugaison complexe les échange.}.

\emph{Cartes régulières et quotients diédraux (page 111).} Le polygone régulier
à $n$ côtés, dessiné sur l'équateur de la sphère, troué en ses sommets, ses
milieux d'arêtes et ses deux centres de faces, donne une surface
$\Sigma_n^{**}$ et une tour
\[
\Sigma_n^{**} \xrightarrow{\ \deg n\ } \Sigma_1^{**} = \mathbb{P}^1 \setminus \{0, \infty, 1, -1\} \xrightarrow{\ \deg 2\ } U_{0,3},
\]
la première flèche étant $z \mapsto z^{n}$, la seconde le quotient par
$z \mapsto 1/z$ ; le composé est un revêtement galoisien de $U_{0,3}$, de
groupe le groupe diédral $D_n^{+}$ des rotations de la sphère qui préservent la
carte\footnote{Les exposants de $\Sigma$, petites piles de $*$ et de points
qui disent quels points sont ôtés, sont lus sans certitude sur leur
disposition, et les diagrammes de quotients de la page ne sont pas
reconstruits. On n'énonce que la tour, vérifiée : $z \mapsto z^n$ puis
$z \mapsto z^{n} + z^{-n}$ à une homographie près, de groupe engendré par
$z \mapsto e^{2i\pi/n} z$ et $z \mapsto 1/z$, d'ordre $2n$.}. La page 112
dessine les douze points $\pm R_i^{\omega}$, $\pm R_i^{\omega'}$ et les
douze valeurs de $P_i^{\omega,\zeta}$.

\pagerange{113}{116}
\subsection*{Le groupoïde fondamental d'une carte trouée (pages 113 à 116)}

Soit $\mathcal{X}$ une carte finie, non nécessairement orientée, $\mathcal{R}$
l'ensemble de ses \emph{repères} (drapeaux), $S = S_0 \sqcup S_1 \sqcup
S_\infty$ l'ensemble de ses sommets, arêtes et faces, identifiés à leurs
centres $R_\varphi$. Un repère $r = (s, a, f)$ s'identifie au triangle de la
subdivision barycentrique de sommets $R_s, R_a, R_f$, d'indices $0, 1,
\infty$, orienté dans le sens $0, 1, \infty$. Les repères adjacents
$\tau_0(r)$, $\tau_1(r)$, $\tau_\infty(r)$ sont ceux qui partagent avec $r$ le
côté opposé au sommet d'indice correspondant, et
\[
(1)\text{–}(5)\qquad \tau_0^2 = \tau_1^2 = \tau_\infty^2 = (\tau_0\tau_\infty)^2 = 1,
\]
\[
\rho_\infty = \tau_1\tau_0,\ \rho_0 = \tau_\infty\tau_1,\ \rho_1 = \tau_0\tau_\infty, \qquad \rho_\infty\rho_1\rho_0 = 1 ,
\]
$\rho_\infty$ tournant d'un cran autour de la face, $\rho_0$ autour du sommet,
$\rho_1$ autour du milieu de l'arête. C'est le groupe cartographique de la page
9 du dossier 141, $\tau_0, \tau_1, \tau_\infty$ y étant notés $\sigma_0,
\sigma_1, \sigma_2$\footnote{Vérifié : $\rho_\infty\rho_1\rho_0 =
\tau_1\tau_0\,\tau_0\tau_\infty\,\tau_\infty\tau_1 = 1$.}.

Sur chaque côté du triangle $r$, près de chacune de ses extrémités, on place un
sommet : $R^{i,j}_r$ est voisin du sommet d'indice $i$, sur le côté qui le
joint au sommet d'indice $j$ (6), avec les seules identifications
$R^{i,j}_r = R^{i,j}_{\tau_k(r)}$, $k \notin \{i, j\}$ (7) — les deux
triangles partagent ce côté. Les arcs sont
\[
(8)\qquad \lambda^i_r : R^{i,i+1}_r \longrightarrow R^{i,i+2}_r, \qquad
(9)\qquad a^i_r : R^{i+1,i+2}_r \longrightarrow R^{i+2,i+1}_r ,
\]
où $i + 1$, $i + 2$ sont les successeurs de $i$ dans l'ordre circulaire $(0,
1, \infty)$ : $\lambda^i_r$ est la portion, intérieure à $r$, du petit cercle
autour du sommet d'indice $i$, et $a^i_r$ longe le côté opposé à ce sommet, de
sorte que $a^i_r = a^i_{\tau_i(r)}$ (10)\footnote{La page écrit (9) $a^i_r :
R^{i+1,i}_r \to R^{i+1,i+2}_r$ (lecture surchargée), qui serait un arc autour
du sommet d'indice $i + 1$, comme $\lambda^{i+1}$, et rendrait (10) et (12)
incohérentes. L'arc le long du côté opposé au sommet $i$ est ce qu'exige
(10), et ce que dit le texte : les $a^i_r$ « entourent » le milieu du repère
« en formant un contour fermé dans le sens défini par l'orientation de $r$ ».}.
Le triangle tronqué est un hexagone, dont le bord donne la relation
\[
(12)\qquad \underbrace{(\lambda^0_r)^{-1} a^1_r}_{b^\infty_r}\ \underbrace{(\lambda^\infty_r)^{-1} a^0_r}_{b^1_r}\ \underbrace{(\lambda^1_r)^{-1} a^\infty_r}_{b^0_r} = 1 ,
\]
lacet basé en $R^{0,1}_r$\footnote{La page dit « lacet en $R^{1\infty}_r$ ».
Avec les arcs (8) et (9) ci-dessus et la composition de droite à gauche, le
produit part de $R^{0,1}_r$ et y revient après avoir fait le tour de
l'hexagone dans le sens de $r$ ; vérifié pas à pas. Le point base d'une
relation est sans effet sur le groupoïde présenté.}.

\emph{Les sommets $R^{ij}_r$, les flèches $\lambda^i_r$ et $a^i_r$, avec les
relations (7), (10), (12), présentent le groupoïde fondamental
$\Pi_1(\widetilde{X}^{**}, \{R^{ij}_r\})$}, où $\widetilde{X}^{**}$ est la
surface à bord obtenue en éclatant réellement la carte en tous les points de
$S$, d'intérieur $X^{**} = X \setminus \{R_\varphi\}$\footnote{Les points et
chemins sont « parfaitement canoniques, pas seulement des classes
d'homotopie », dans la réalisation conforme canonique de la carte. Que ce soit
une présentation est l'énoncé de van Kampen pour le complexe formé du graphe
des $\lambda$ et des $a$ et des hexagones : son complémentaire est une réunion
de disques épointés autour des $R_\varphi$.}. Les notations gardent un sens
pour toute carte triangulaire « partielle » de type $\{0, 1, \infty\}$ ; la
relation $(\tau_0\tau_\infty)^2 = 1$, qui dit que les sommets de type $1$ sont
de valence $4$, est la seule qui la fasse venir de la subdivision barycentrique
d'une carte. On introduit les « rotations d'un double cran »
\[
(13)\qquad \Lambda^i_r = (\lambda^i_{\tau_{i+1}(r)})^{-1}\lambda^i_r : R^{i,i+1}_r \to R^{i,i+1}_{\tau_{i+1}(r)}, \qquad
\Lambda'^i_r = \lambda^i_{\tau_{i-1}(r)}(\lambda^i_r)^{-1} : R^{i,i-1}_r \to R^{i,i-1}_{\tau_{i-1}(r)},
\]
qui font le tour du sommet d'indice $i$ à travers deux triangles, et
satisfont $\Lambda^i_r\Lambda^i_{\tau_{i+1}(r)} = 1$,
$\Lambda'^i_r\Lambda'^i_{\tau_{i-1}(r)} = 1$ (14) — des relations de
redondance\footnote{Vérifié sur (13) ; les indices $\tau_{i\pm 1}$ sont lus
avec hésitation.}.

\pagerange{117}{120}
\subsection*{Présentations réduites (pages 117 à 120)}

En éliminant les sommets sur les côtés de type $1$ (joignant un sommet et un
centre de face) et les arcs $a^1_r$ qui les joignent, on garde les $R^{ij}_r$
avec $1 \in \{i, j\}$, les arcs $\lambda^1_r$, $\Lambda^0_r$,
$\Lambda^\infty_r$, $a^0_r$, $a^\infty_r$, les relations de redondance (14), et
une relation par quadrilatère $r \cup \tau_1(r)$ :
\[
(16)\qquad \Lambda'^\infty_{\tau_1(r)}\,(\beta^1_{\tau_1(r)})^{-1}\,\Lambda^0_r\,(\beta^1_r)^{-1} = 1, \qquad \beta^1_r = a^0_r(\lambda^1_r)^{-1}a^\infty_r ,
\]
stable par $r \mapsto \tau_1(r)$, de sorte qu'on peut se borner, si $X$ est
orientable, aux repères d'une orientation\footnote{Une première version de
(16) est barrée. L'exposant de la première parenthèse est surchargé, lu $-1$
avec hésitation ; on rapporte la relation sans l'avoir reconstruite pas à
pas. La marge propose de voir le résultat comme une décomposition en
« losanges », les sommets de type $1$ étant arrondis.}.

En ne gardant que les arcs $a^\infty_r$ et les sommets $R^{0,1}_r$,
$R^{1,0}_r$ qu'ils joignent, avec les doubles crans autour des sommets de type
$0$ et $1$,
\[
(18)\qquad a^\infty_r : R^{0,1}_r \to R^{1,0}_r, \qquad \Lambda^0_r : R^{0,1}_r \to R^{0,1}_{\tau_1(r)}, \qquad \Lambda'^1_r : R^{1,0}_r \to R^{1,0}_{\tau_0(r)},
\]
les seules relations sont de redondance (19) : \emph{le groupoïde est libre}.
En effet le graphe correspondant — de petits cercles autour des sommets de type
$0$ et $1$, reliés le long des demi-arêtes — a pour complémentaire des disques
épointés, et la surface se rétracte sur lui. Cela fait trois présentations, sur
$S \setminus S_i$, $i = \infty, 0, 1$. Enfin, en ne gardant que les
$R^{0,1}_r$, on obtient un groupoïde libre engendré par
\[
(21)\qquad A^{0,1}_r = (a^\infty_{\tau_0(r)})^{-1}\,\Lambda'^1_r\,a^\infty_r : R^{0,1}_r \to R^{0,1}_{\tau_0(r)}, \qquad \Lambda^0_r : R^{0,1}_r \to R^{0,1}_{\tau_1(r)},
\]
avec les seules relations de redondance $\Lambda^0_r\Lambda^0_{\tau_1(r)} = 1$,
$A^{0,1}_r A^{0,1}_{\tau_0(r)} = 1$ (22)\footnote{La page écrit
$(\Lambda'^1_r)^{-1}$ au milieu de $A^{0,1}_r$ (définition surchargée, lue
avec hésitation). Pour que le chemin se compose, il faut aller de $R^{1,0}_r$
à $R^{1,0}_{\tau_0(r)}$, ce que fait $\Lambda'^1_r$ ; son inverse va dans
l'autre sens.}. Au total, $10 = 1 + 3 + 3 + 3$ présentations, selon les
ensembles de sommets retenus.

\pagerange{120}{121}
\subsection*{Le bouchage des trous (pages 120 et 121)}

« Dans chacun des cas envisagés, le “bouchage” d'un trou $R_\varphi$ se fait de
façon essentiellement triviale » : il suffit d'exprimer, en termes des flèches
génératrices choisies, le « lacet » autour de $R_\varphi$ — pour un sommet de
type $0$, le produit des doubles crans $\Lambda^0$ qui en font le tour —,
« puis d'exprimer que ce lacet est $\simeq 1$ »\footnote{La phrase s'arrête au
bas de la page 120 et reprend en tête de la page 121, sa page 7, dont le reste
est blanc ; l'exemple du produit des $\Lambda^0$ est de l'édition. Les
« $R_s$ » de la page 121 sont les centres $R_\varphi$ de type sommet.}. Le
groupoïde fondamental de la carte dont on a bouché ce trou est le quotient par
cette relation, par le théorème de van Kampen.

C'est la réponse, au niveau combinatoire et pour une carte donnée, à la
difficulté laissée en suspens aux pages 17 à 20 : là, le bouchage des trous ne
trouvait pas de place, à côté des revêtements, dans une catégorie dont on
pût interpréter la complétion profinie ; ici, sur des présentations
explicites, il n'est qu'une relation de plus. Le dossier s'arrête sur ce
constat.

\end{document}
